\section*{Chapitre \ref{ch:iv3:algorithms} --- Algorithmes}
\markboth{\protect\omTocMark{}{Réponses~: Algorithmes}}{\protect\omTocMark{}{Réponses~: Algorithmes}}

\begin{solution}{iq:iv3:algorithms:1}
On garde le plus haut prix vu jusqu'ici, et à chaque mise à jour du prix on retient le plus grand écart en dessous de lui, comme dans le \cref{lst:iv3:al:drawdown}. Le temps est en $O(n)$ et la mémoire en $O(1)$, et l'exemple renvoie 40~: la chute de 120 à 80.
\omcode{code/interviews-3/01-algorithms/python/iv_drawdown.py}{1}{8}{Un seul parcours avec un pic courant.}\label{lst:iv3:al:drawdown}

Un prix ultérieur ne peut être mesuré que par rapport à un pic qui le précède, et le pic courant est le meilleur de ces pics à chaque instant~; aucune paire antérieure n'a donc besoin d'être réexaminée. Une liste sans aucune chute, ou une liste vide, renvoie 0.

\iqwrong{«~Prix le plus haut moins prix le plus bas.~» Cela ignore l'ordre. Dans l'exemple, cela donne 130 moins 80, soit 50, mais 130 vient après 80, donc ce mouvement est une hausse. Pour \texttt{[1, 2, 3, 10]}, cela donne 9, alors que le vrai drawdown vaut 0.}
\iqfollow{Et en pourcentage du pic~?}{Le même parcours, avec \texttt{(peak - p) / peak}. Il peut retenir une autre chute~: \texttt{[100, 50, 1000, 600]} a sa pire chute en valeur à 400, depuis 1000, et sa pire chute en pourcentage à 50~\%, depuis 100.}
\iqfollow{Le gain maximal d'un achat suivi d'une vente ultérieure~?}{L'image miroir~: on garde le plus bas courant et on met à jour le plus grand écart au-dessus de lui.}
\end{solution}

\begin{solution}{iq:iv3:algorithms:2}
On parcourt la liste une fois avec un dictionnaire qui associe chaque valeur à sa position, et on cherche le partenaire avant d'insérer l'élément courant, comme dans le \cref{lst:iv3:al:twosum}. Le temps moyen est en $O(n)$ et la mémoire en $O(n)$. Pour \texttt{[3, 4]} avec la cible 6, aucune paire n'est signalée.
\omcode{code/interviews-3/01-algorithms/python/iv_two_sum.py}{1}{8}{Chercher, puis insérer.}\label{lst:iv3:al:twosum}

Chaque paire est trouvée quand son élément le plus tardif arrive, car son élément antérieur figure alors dans le dictionnaire~; un seul parcours suffit. L'ordre des deux lignes dans la boucle fait tout l'intérêt~: comme l'élément courant n'est pas encore inséré, il ne peut pas se former en paire avec lui-même, alors que \texttt{[3, 3]} avec la cible 6 donne toujours les positions 0 et 1.

\iqwrong{Remplir d'abord le dictionnaire, puis chercher. Sur \texttt{[3, 4]} avec la cible 6, on trouve 3 comme partenaire de 3 et l'on renvoie les positions 0 et 0.}
\iqfollow{Et si la liste est déjà triée~?}{On utilise deux pointeurs partant des extrémités~: on déplace le gauche vers le haut si la somme est trop petite, et le droit vers le bas si elle est trop grande. Cela demande une mémoire en $O(1)$, alors que trier d'abord coûte $O(n \log n)$ et perd les positions, sauf si l'on trie aussi les positions avec les valeurs.}
\iqfollow{Quel est le pire cas du dictionnaire~?}{Si beaucoup de clés entrent en collision, une recherche se dégrade en parcours, et toute la boucle vers $O(n^2)$. Sur une entrée non fiable, les tables ordonnées donnent une borne ferme en $O(n \log n)$.}
\end{solution}

\begin{solution}{iq:iv3:algorithms:3}
Les exécutions sont 300 à 100, 200 à 101 et 100 à 103, pour un total de 60\,500 et un prix moyen d'environ 100.83. Le \Cref{lst:iv3:al:sweep} parcourt les niveaux. Le temps est proportionnel au nombre de niveaux touchés, et non au nombre d'actions.
\omcode{code/interviews-3/01-algorithms/python/iv_sweep.py}{1}{11}{Parcours du carnet d'ordres.}\label{lst:iv3:al:sweep}

Le coût vaut $300 \times 100 + 200 \times 101 + 100 \times 103 = 60\,500$, et $60\,500 / 600 \approx 100.83$, soit 0.83 au-dessus du meilleur cours vendeur~: cet écart est le prix de la taille de l'ordre. Les prix et les tailles restent entiers, en ticks et en lots, si bien que le coût cumulé est exact.

L'écueil que la fonction ne doit pas masquer est un ordre plus grand que le carnet. Pour 1200 actions, elle exécute 1000 et renvoie 200 non exécutées, pour un coût de 101\,700. Le prix moyen par action exécutée vaut 101.70, et diviser par les 1200 demandées donnerait un faux 84.75. Il faut donc calculer la moyenne sur ce qui a été exécuté, et signaler ce qui ne l'a pas été~; la fonction laisse aussi le carnet tel qu'elle l'a trouvé, ce qu'un appelant peut vouloir ou non, et il faut préciser lequel on a retenu.

\iqfollow{Combien peut-on acheter à un prix limite de 101 ou en dessous~?}{On s'arrête au premier niveau au-dessus de la limite~: ici, 500 actions. Pour de nombreuses requêtes de ce type sur un même instantané, on stocke les tailles cumulées et l'on utilise une recherche dichotomique.}
\end{solution}

\begin{solution}{iq:iv3:algorithms:4}
On garde un tampon circulaire des $k$ dernières valeurs et un total courant, comme dans le \cref{lst:iv3:al:rolling}~: à chaque mise à jour, on retranche la valeur qui sort et l'on ajoute celle qui arrive. Avec ces flottants, le dernier appel renvoie $1/3$, soit environ 0.333, alors que la vraie moyenne vaut 1.
\omcode{code/interviews-3/01-algorithms/python/iv_rolling_mean.py}{1}{15}{Un tampon circulaire et un total courant.}\label{lst:iv3:al:rolling}

C'est l'arrondi qui casse le calcul. Les doubles au voisinage de $10^{16}$ sont espacés de 2, si bien que \texttt{1e16 + 1.0} s'arrondit à \texttt{1e16}, et les deux valeurs 1.0 disparaissent du total. Quand \texttt{1e16} quitte la fenêtre, le total tombe à 0, puis le dernier 1.0 est ajouté, de sorte que la moyenne vaut $1/3$ sur une fenêtre qui contient trois valeurs 1.0. L'erreur reste ensuite dans le total.

La même dérive se produit sur des prix ordinaires, mais plus lentement~: après deux millions de mises à jour avec des prix voisins de 100, le total de notre test n'égale plus la somme de ses dix valeurs, bien qu'il s'en écarte de beaucoup moins d'un millionième. Deux remèdes~: tenir les prix en ticks entiers, auquel cas la somme reste exacte en Python, et en C++ jusqu'au débordement~; ou recalculer périodiquement le total à partir du tampon.

\iqwrong{«~Sommer le tampon à chaque appel.~» C'est correct, mais cela coûte $O(k)$ par mise à jour, ce que la question exclut.}
\end{solution}

\begin{solution}{iq:iv3:algorithms:5}
On utilise une pile~: on empile chaque parenthèse ouvrante, et à chaque parenthèse fermante on dépile et l'on compare, comme dans le \cref{lst:iv3:al:brackets}. Le temps est en $O(n)$ et la mémoire au plus en $O(n)$. \texttt{([)]} renvoie faux, et la chaîne vide renvoie vrai.
\omcode{code/interviews-3/01-algorithms/python/iv_brackets.py}{1}{10}{Une pile de parenthèses ouvertes.}\label{lst:iv3:al:brackets}

Il y a deux façons d'échouer, et les deux doivent être détectées~: une parenthèse fermante avec une pile vide ou un mauvais partenaire, et une pile restée non vide à la fin, comme pour \texttt{((}.

\iqwrong{Un compteur par type de parenthèse. Sur \texttt{([)]}, chaque compteur monte de un puis redescend à zéro sans jamais devenir négatif, si bien qu'il renvoie vrai. L'ordre est ce qu'un compteur ne voit pas.}
\iqfollow{Avec seulement des parenthèses rondes, peut-on économiser la mémoire~?}{Oui~: un seul compteur suffit, puisqu'il n'y a qu'un seul partenaire possible. Il ne doit jamais devenir négatif, et il doit finir à 0.}
\iqfollow{Quelle est la plus longue portion valide d'une chaîne de parenthèses rondes~?}{On garde une pile de positions, avec la position juste avant la portion courante au fond. Chaque fermeture dépile, et la longueur vaut la position courante moins le nouveau sommet. Pour \texttt{)()())}, elle vaut 4.}
\end{solution}

\begin{solution}{iq:iv3:algorithms:6}
Elle renvoie 1 pour $x = 2$ et 5 pour $x = 4$. Le \Cref{lst:iv3:al:lowerbound} maintient l'invariant selon lequel la réponse se situe entre \texttt{lo} et \texttt{hi}, bornes comprises~; c'est pourquoi \texttt{hi} part de $n$ et non de $n - 1$~: $n$ lui-même est une réponse possible.
\omcode{code/interviews-3/01-algorithms/python/iv_lower_bound.py}{1}{9}{Première position au moins égale à la cible.}\label{lst:iv3:al:lowerbound}

Si la valeur du milieu est trop petite, la réponse est plus à droite, au-delà du milieu. Sinon, le milieu lui-même peut être la réponse, si bien que \texttt{hi} se déplace jusqu'à lui et non au-delà. L'intervalle se réduit à chaque étape, car \texttt{mid} est toujours strictement inférieur à \texttt{hi}, et la boucle se termine.

En C++, \texttt{lo + hi} peut dépasser le plus grand \texttt{int} dès que les positions dépassent environ un milliard, et un débordement signé est un comportement indéfini. \texttt{lo + (hi - lo) / 2} ne peut pas déborder.

\iqwrong{S'arrêter dès que \texttt{xs[mid] == x}. On trouve alors une position de la valeur, et non la première~: sur cinq copies de la valeur, on s'arrête tout de suite au milieu.}
\iqfollow{Combien de fois $x$ apparaît-il~?}{La première position au-dessus de $x$ moins la première position au moins égale à $x$~: 4 moins 1 vaut 3 pour la liste et $x = 2$. Chacune est une recherche dichotomique, si bien que le comptage coûte $O(\log n)$ quel que soit le nombre de copies.}
\end{solution}

\begin{solution}{iq:iv3:algorithms:7}
Après A, B, A, C, B, elle renvoie C. Le \Cref{lst:iv3:al:firstunique} garde les symboles vus une seule fois dans un \texttt{OrderedDict}, et ceux vus plus souvent dans un ensemble.
\omcode{code/interviews-3/01-algorithms/python/iv_first_unique.py}{1}{19}{Un dictionnaire ordonné et un ensemble.}\label{lst:iv3:al:firstunique}

Un symbole passe au plus une fois du dictionnaire à l'ensemble, et jamais en sens inverse. Un \texttt{OrderedDict} est une table de hachage enfilée sur une liste chaînée, si bien que supprimer au milieu et lire la première clé, qui est la réponse, coûtent tous deux un temps constant en moyenne. La mémoire est proportionnelle au nombre de symboles distincts.

En C++, la même structure est une \texttt{std::list} de candidats et une \texttt{unordered\_map} qui associe chaque symbole à sa place dans la liste, une place nulle signifiant répété.

\iqwrong{Un \texttt{dict} ordinaire, puisqu'il conserve lui aussi l'ordre d'insertion. En CPython, \texttt{next(iter(d))} saute les entrées supprimées au début de sa table, et celles-ci restent là jusqu'à ce qu'une insertion la redimensionne. Quand les symboles les plus anciens ne cessent de sortir, chaque appel coûte $O(n)$, et le flux devient quadratique.}
\iqfollow{Une file des arrivées à la place~?}{Cela fonctionne~: on retire définitivement l'élément de tête tant qu'il est répété, et on lit ce qui reste. Chaque symbole n'est retiré qu'une fois, si bien que le coût est constant amorti, ce que l'énoncé permet. Un parcours depuis le début à chaque appel est la version en $O(n)$.}
\iqfollow{Et si vous ne le voulez qu'à la fin, et non après chaque symbole~?}{On compte tout, puis on reparcourt l'entrée une fois~: c'est plus simple, à condition que l'entrée soit encore disponible.}
\end{solution}

\begin{solution}{iq:iv3:algorithms:8}
On parcourt la liste avec deux pointeurs et l'on retourne chaque lien, comme dans le \cref{lst:iv3:al:revlist}. Le temps est en $O(n)$ et la mémoire supplémentaire en $O(1)$, et une liste vide renvoie une liste vide, car la boucle ne s'exécute pas et \texttt{prev} vaut toujours \texttt{None}.
\enlargethispage{-3\baselineskip}%
\omcode{code/interviews-3/01-algorithms/python/iv_reverse_list.py}{1}{10}{Retourner chaque lien.}\label{lst:iv3:al:revlist}

À chaque étape, tout ce qui précède \texttt{head} est déjà retourné, avec \texttt{prev} en tête. L'affectation de trois noms est évaluée d'abord à droite, si bien qu'aucun pointeur n'est perdu. Écrite en C++, il faut un temporaire pour \texttt{next}, et l'oublier est le bogue habituel.

La version récursive tient en trois lignes, mais elle utilise un cadre de pile par nœud. La limite de récursion par défaut de Python est 1000, si bien que sur 100\,000 nœuds elle lève \texttt{RecursionError}. En C++, le même code peut faire déborder la pile, et ce plantage n'a pas d'exception à attraper.

\iqwrong{Inverser les valeurs sur place en les copiant dans une liste puis en les recopiant. Cela répond à une autre question~: un nœud tenu ailleurs a désormais une autre valeur, et la méthode utilise une mémoire en $O(n)$.}
\iqfollow{Inverser seulement de la position $a$ à la position $b$~?}{On va jusqu'au nœud qui précède $a$, on retourne $b - a + 1$ liens de la même façon, puis on recolle les deux extrémités.}
\end{solution}

\begin{solution}{iq:iv3:algorithms:9}
Elle efface le bit de poids faible à 1~: 12 vaut 1100 en binaire, et \texttt{12 \& 11} vaut 1000 en binaire, soit 8. Le \Cref{lst:iv3:al:bits} s'en sert pour les trois tâches.
\omcode{code/interviews-3/01-algorithms/python/iv_bits.py}{1}{14}{Trois usages du bit de poids faible à 1.}\label{lst:iv3:al:bits}

Soustraire 1 transforme le bit de poids faible à 1 en 0 et chaque zéro situé en dessous en 1, et le \texttt{\&} élimine ces bits. Ainsi \texttt{x} est une puissance de deux exactement quand le résultat vaut 0, puisqu'il ne contenait alors qu'un seul bit à 1. Le test doit écarter 0 à part, car \texttt{0 \& -1} vaut aussi 0. Le comptage ne boucle qu'une fois par bit à 1, et non une fois par bit.

L'arrondi fonctionne, car \texttt{x - 1} a la même longueur en bits que l'exposant de la réponse~: 1000 s'arrondit à 1024, et 1024 reste 1024. Un tampon circulaire de taille $2^m$ transforme le retour à zéro d'un compteur croissant \texttt{i} en \texttt{i \& (size - 1)}, qui est une seule instruction peu coûteuse, là où \texttt{\%} sur une taille inconnue est une division entière. Le compteur peut aussi déborder sans risque un type non signé, car le retour à zéro d'une puissance de deux s'aligne sur le masque.

\begin{sloppypar}\iqfollow{Qu'utiliseriez-vous en C++20 en production~?}{L'en-tête \texttt{<bit>} fournit \texttt{std::has\_single\_bit}, \texttt{std::popcount} et \texttt{std::bit\_ceil}, qui se compilent en une seule instruction lorsque le processeur en possède une.}\end{sloppypar}
\end{solution}

\begin{solution}{iq:iv3:algorithms:10}
On fait le XOR de tous les $1, \dots, n$ avec toutes les valeurs présentes, comme dans le \cref{lst:iv3:al:missing}. Chaque nombre présent s'annule, et il reste le nombre manquant. Pour l'exemple, cela renvoie 4. Le temps est en $O(n)$ et la mémoire en $O(1)$.
\omcode{code/interviews-3/01-algorithms/python/iv_missing.py}{1}{7}{Tout ce qui est présent s'annule.}\label{lst:iv3:al:missing}

L'autre classique est la somme~: $n(n + 1)/2$ moins le total de ce qui est présent. Elle déborde d'un entier signé de 64 bits dès que $n(n + 1)/2$ dépasse $2^{63} - 1$, ce qui arrive vers $n = 4.3 \times 10^9$. Les numéros de séquence peuvent atteindre cette valeur sur un flux de longue durée, et le XOR ne peut jamais déborder, ce qui en fait le plus sûr des deux.

Si le fichier est en ordre, aucune des deux n'est nécessaire~: il suffit de comparer chaque nombre au précédent plus un.

\iqwrong{Trier d'abord, ou utiliser un ensemble. Les deux donnent le bon nombre, et aucune ne respecte la limite de mémoire.}
\iqfollow{Et si deux nombres manquent~?}{Le XOR vaut maintenant $a \oplus b$. Il est non nul, donc on prend un de ses bits à 1, bit sur lequel $a$ et $b$ diffèrent. Puis on refait un XOR, sur les seuls nombres dont ce bit vaut 1, ce qui donne $a$, et $b$ s'en déduit. Pour \texttt{[3, 5, 2]} avec $n = 5$, la paire manquante est 1 et 4.}
\end{solution}

\begin{solution}{iq:iv3:algorithms:11}
On conserve une seconde pile du minimum à chaque profondeur où il a changé, comme dans le \cref{lst:iv3:al:minstack}. Les trois opérations sont en $O(1)$ dans le pire cas, et la mémoire supplémentaire vaut au plus $O(n)$.
\omcode{code/interviews-3/01-algorithms/python/iv_min_stack.py}{1}{18}{Une seconde pile des minima.}\label{lst:iv3:al:minstack}

Une seule variable échoue sur \texttt{pop}~: lorsque le minimum courant sort, la variable ne permet pas de retrouver le précédent sans parcourir la pile. La seconde pile contient précisément cet historique.

La comparaison dans \texttt{push} doit être \texttt{<=} et non \texttt{<}. Empiler 1 et 1 puis dépiler une fois~: le minimum vaut toujours 1, et avec \texttt{<} l'unique entrée de la seconde pile aurait été dépilée, ne laissant rien.

\iqwrong{Un tas à côté de la pile. Il trouve le minimum, mais dépiler devrait supprimer une valeur quelconque du tas, ce qui coûte davantage et n'apporte rien de plus ici.}
\iqfollow{Même question pour une file~?}{On la construit à partir de deux piles de ce type, de sorte que le minimum soit le plus petit des deux minima des piles, en $O(1)$ amorti~; ou on utilise une deque de valeurs croissantes, comme pour un minimum sur une fenêtre glissante.}
\end{solution}

\begin{solution}{iq:iv3:algorithms:12}
On empile sur \texttt{inbox}. Pour dépiler, on prend dans \texttt{outbox}, et lorsque celle-ci est vide, on transfère d'abord tout \texttt{inbox}, ce qui l'inverse dans l'ordre d'arrivée, comme dans le \cref{lst:iv3:al:twostack}. \index{analyse amortie}Chaque élément n'est déplacé qu'une fois au plus, donc $n$ opérations font au plus $n$ déplacements au total~: $O(1)$ amorti.
\omcode{code/interviews-3/01-algorithms/python/iv_two_stack_queue.py}{1}{14}{Transférer seulement lorsque la sortie est vide.\protect\hfil\protect\break}\label{lst:iv3:al:twostack}

Suivons un seul élément~: empilé une fois, déplacé une fois, dépilé une fois. Cela fait trois opérations, quelle que soit la longueur de la file. Le dépilement coûteux paie pour les opérations peu coûteuses qui l'ont précédé, et c'est ce que signifie «~amorti~». Il ne s'agit pas d'une moyenne sur des entrées aléatoires~: la propriété vaut pour toute séquence.

L'ordre est préservé, car \texttt{outbox} n'est jamais rechargé tant qu'il n'est pas vide, de sorte que tout ce qu'il contient est plus ancien que tout ce qui se trouve dans \texttt{inbox}.

Sur un chemin soumis à une limite de latence, on ne l'utiliserait pas. Le $O(1)$ amorti autorise tout de même un \texttt{pop} qui déplace un million d'éléments, et c'est précisément ce message qui dépasse la limite. Un tampon circulaire (ring buffer) donne une véritable constante dans le pire des cas.

\iqwrong{«~Remettre tout dans \texttt{inbox} après chaque dépilement.~» Le résultat reste correct, mais coûte $O(n)$ par opération.}
\end{solution}

\begin{solution}{iq:iv3:algorithms:13}
Le \Cref{lst:iv3:al:ticks} découpe le texte au point et additionne des entiers. La version en flottants donne 114 pour le texte 1.15, soit une erreur d'un tick. Le texte -0.05 doit donner $-5$.
\omcode{code/interviews-3/01-algorithms/python/iv_ticks.py}{1}{8}{Du texte à des ticks entiers, sans flottant.}\label{lst:iv3:al:ticks}

1.15 n'a pas de forme binaire exacte. Le double le plus proche est légèrement inférieur, le produit vaut 114.99999999999999, et \texttt{int} tronque vers zéro. Arrondir d'abord, avec \texttt{round(float(s) * 100)}, corrige cet exemple, mais la valeur passe encore par un flottant, ce qui échoue pour des prix assez grands pour que le chiffre se perde, et cela ne permet pas de refuser un texte comportant trop de décimales.

Le signe doit d'abord être retiré, puis rétabli à la fin. Un analyseur qui lit la partie entière comme un entier voit \texttt{-0} pour -0.05. C'est simplement 0, donc le signe disparaît et le résultat vaut $+5$~: un cours acheteur situé du mauvais côté de zéro. Cela compte sur les fourchettes, qui peuvent être négatives.

\iqwrong{«~Comparer le flottant à un epsilon.~» Le code de trading manipule alors les prix comme des clés non exactes, et deux flux peuvent ne pas s'accorder sur le fait qu'un prix soit égal à un niveau de prix.}
\iqfollow{Et en C++~?}{\texttt{std::from\_chars} analyse chaque partie entière sans toucher à la locale ni allouer de mémoire, ce qui convient à un chemin critique (hot path).}
\end{solution}

\begin{solution}{iq:iv3:algorithms:14}
Le code renvoie une somme de 6, des positions 1 à 3, pour la première liste, et $-1$, l'unique élément en position 1, pour la seconde. Le \Cref{lst:iv3:al:kadane} est l'\index{algorithme de Kadane}algorithme de Kadane~: $O(n)$ en temps et $O(1)$ en mémoire.
\omcode{code/interviews-3/01-algorithms/python/iv_kadane.py}{1}{11}{Le meilleur segment se terminant ici.}\label{lst:iv3:al:kadane}

\texttt{here} est la meilleure somme sur les segments se terminant au jour courant. Si le meilleur segment se terminant la veille est négatif, l'abandonner fait mieux, et le nouveau segment commence donc aujourd'hui. Sinon, le prolonger fait mieux. Tout segment se termine quelque part, donc la plus grande valeur jamais atteinte par \texttt{here} est la réponse.

L'énoncé dit non vide. C'est pourquoi le code part du premier élément et non de 0~: partir de 0 renvoie 0 sur une liste entièrement négative, qui est la somme d'aucun jour. Une liste vide n'a aucun segment de ce type, si bien que le code lève \texttt{IndexError}~: il faut préciser que c'est un choix.

\iqwrong{Trier les jours, ou additionner les jours positifs. Les deux ignorent que les jours doivent être contigus~: les jours positifs de la première liste font 7, et aucun segment n'atteint cette somme.}
\iqfollow{Quel est le lien avec le drawdown maximal~?}{On prend le profit cumulé. La somme d'un segment est une valeur cumulée moins une valeur antérieure, donc il s'agit de la plus forte hausse, et le drawdown maximal se calcule par le même parcours, avec le signe opposé.}
\end{solution}

\begin{solution}{iq:iv3:algorithms:15}
Le code renvoie \texttt{[0, 0, 24, 0]}. Le \Cref{lst:iv3:al:product} effectue deux passes~: la première multiplie tout ce qui se trouve à gauche de chaque position, la seconde tout ce qui se trouve à droite. Cela demande un temps en $O(n)$, et la liste de sortie est le seul tableau.
\omcode{code/interviews-3/01-algorithms/python/iv_product_except.py}{1}{11}{Une passe à gauche et une passe à droite.}\label{lst:iv3:al:product}

Le produit total de cette liste vaut 0, donc le diviser par 0 à la position 2 échoue, et le diviser par 2 à la position 0 donne 0, ce qui n'est juste que par hasard pour les autres positions. Avec une division, il faudrait compter les zéros~: deux zéros ou plus rendent toutes les sorties nulles, et un seul zéro rend nulles toutes les sorties sauf celle de cette position.

La liste \texttt{[2, 3, 0, 0]} illustre le cas de deux zéros, et elle renvoie uniquement des zéros.

\iqwrong{Utiliser des logarithmes et soustraire. Les zéros et les valeurs négatives n'ont pas de logarithme, et l'arrondi donne de toute façon des résultats inexacts.}
\iqfollow{Avec des entiers 64 bits en C++, qu'est-ce qui peut mal tourner~?}{Les produits cumulés peuvent déborder même lorsque chaque sortie tient, par exemple si un zéro vient après de grandes valeurs, et un débordement signé est un comportement indéfini. Les entiers de Python ne débordent pas.}
\end{solution}

\begin{solution}{iq:iv3:algorithms:16}
Python 3 donne $-4$ et $1$, car \texttt{//} arrondit vers le bas et le reste prend le signe du diviseur. C++20 donne $-3$ et $-1$, car \texttt{/} tronque vers zéro et le reste prend le signe du dividende. En C++, le \cref{lst:iv3:al:floorsnap} corrige cela.
\omcode{code/interviews-3/01-algorithms/cpp/floor_snap.cpp}{5}{8}{Arrondi vers le bas en C++.}\label{lst:iv3:al:floorsnap}

Dans les deux langages, le quotient multiplié par le diviseur, plus le reste, redonne la valeur, ce qui explique que les deux réponses se correspondent~: $-4 \times 2 + 1$ et $-3 \times 2 - 1$ valent toutes deux $-7$.

En Python, la fonction tient en une seule expression, \texttt{p - p \% tick}, car le reste n'est jamais négatif pour un tick positif. Ainsi $-7$ avec un tick de 5 donne $-7 - 3 = -10$. Le reste C++ pour ces valeurs vaut $-2$, et ne soustraire que ce reste donne $-5$, ce qui va dans le mauvais sens~: on arrondit vers zéro, et non vers le bas. Donc, lorsque le reste est négatif, la fonction retire un tick de plus.

\iqwrong{\texttt{x / tick * tick} en C++. Il donne $-5$ pour $-7$, dans le mauvais sens sur une fourchette négative, et il paraît juste sur chaque test positif.}
\iqfollow{Ce comportement de C++ est-il garanti~?}{Oui~: la norme indique que la division entière supprime la partie fractionnaire, et que $(a/b) \times b + a \% b = a$. Les seules exceptions sont la division par zéro et la division de la valeur la plus négative par $-1$, qui est un comportement indéfini.}
\end{solution}

\begin{solution}{iq:iv3:algorithms:17}
On garde une deque des positions dont les valeurs décroissent strictement, comme dans le \cref{lst:iv3:al:winmax}. La tête de la deque est toujours le maximum de la fenêtre. L'exemple renvoie \texttt{[3, 5, 5, 5]}.
\omcode{code/interviews-3/01-algorithms/python/iv_window_max.py}{4}{15}{Une deque monotone.}\label{lst:iv3:al:winmax}

Quand \texttt{x} arrive, toute valeur plus ancienne qui n'est pas plus grande ne pourra plus jamais être un maximum~: \texttt{x} est au moins aussi grande et quitte la fenêtre plus tard. Ces valeurs sont donc retirées par l'arrière. Ensuite, la tête est retirée si elle est sortie de la fenêtre glissante. Ce qui reste est décroissant, et sa première valeur est la plus grande de la fenêtre.

Une seule étape peut retirer jusqu'à $k$ éléments, ce qui explique que la borne par étape ressemble à $k$. Le bon décompte se fait sur toute l'exécution~: chaque position est ajoutée une fois et retirée au plus une fois, donc l'ensemble des boucles \texttt{while} effectue au plus $n$ retraits. Le total vaut $O(n)$, amorti, quel que soit $k$, avec une mémoire en $O(k)$.

\iqwrong{Un tas sur la fenêtre. Cela fonctionne avec une suppression paresseuse, mais les entrées périmées s'accumulent~: sur une liste croissante, aucune n'atteint la racine, donc le coût est en $O(n \log n)$ et la mémoire en $O(n)$. Même $O(n \log k)$ serait exclu par l'énoncé.}
\iqfollow{Et le plus grand moins le plus petit dans chaque fenêtre~?}{On maintient deux deques, une pour le maximum et une pour le minimum, et on soustrait.}
\iqfollow{Et une fenêtre des 60 dernières secondes d'un flux, plutôt que des $k$ derniers éléments~?}{La même deque, qui contient aussi les horodatages~: on retire de la tête les éléments trop anciens selon le temps. La mémoire est alors bornée par le nombre d'événements à l'intérieur de l'intervalle de temps, et non par une constante que l'on a choisie.}
\end{solution}

\begin{solution}{iq:iv3:algorithms:18}
On utilise des sommes préfixes et une table de hachage, comme dans le \cref{lst:iv3:al:subsum}. Un segment allant de $i$ à $j$ a pour somme $T$ exactement lorsque $\text{préfixe}_j - \text{préfixe}_{i} = T$, donc à chaque $j$, on compte combien de préfixes antérieurs valent $\text{préfixe}_j - T$. Le temps moyen est en $O(n)$ et la mémoire en $O(n)$, et l'exemple comporte 3 segments de ce type.
\omcode{code/interviews-3/01-algorithms/python/iv_subarray_sum.py}{1}{11}{Comptage des sommes préfixes déjà rencontrées.}\label{lst:iv3:al:subsum}

Pour \texttt{[1, -1, 1]}, les préfixes valent 0, 1, 0, 1. Les trois segments sont le premier élément, le dernier élément et la liste entière. Le compteur part avec un préfixe vide, de sorte qu'un segment commençant en position 0 soit compté, et le préfixe courant n'est ajouté qu'après la recherche, de sorte qu'avec $T = 0$ un segment vide ne soit pas compté.

\iqwrong{Deux pointeurs. L'agrandissement et le rétrécissement reposent sur une somme qui croît quand la fenêtre grandit et décroît quand elle rétrécit, ce qui exige des valeurs jamais négatives. Avec une valeur négative, une fenêtre au-dessus de $T$ peut revenir à $T$ si on l'agrandit davantage. Sur cet exemple, elle n'en compte que 2 sur 3, puisqu'elle ne garde qu'une fenêtre par extrémité, alors que la dernière position termine deux segments dont la somme vaut 1.}
\iqfollow{Et les segments dont la somme est divisible par $m$~?}{On garde les préfixes modulo $m$ et on compte les valeurs égales déjà rencontrées. Attention au signe~: il faut ramener le reste dans $0, \dots, m-1$, ce que fait automatiquement le \texttt{\%} de Python, mais pas celui de C++.}
\iqfollow{La mémoire pose problème, et toutes les valeurs sont positives. Et maintenant~?}{Dans ce cas, les deux pointeurs conviennent, avec une mémoire supplémentaire en $O(1)$.}
\end{solution}

\begin{solution}{iq:iv3:algorithms:19}
On trie les ouvertures et les fermetures ensemble comme des événements, puis on balaie la liste, comme dans le \cref{lst:iv3:al:alive}. À temps égal, les fermetures doivent venir avant les ouvertures. Le temps est en $O(n \log n)$ et la mémoire en $O(n)$. Pour les deux ordres de l'exemple, la réponse est 1, atteinte pour la première fois au temps 1.
\omcode{code/interviews-3/01-algorithms/python/iv_alive.py}{4}{17}{Un balayage des ouvertures et des fermetures.}\label{lst:iv3:al:alive}

Entre deux événements, le compte ne change pas, donc son maximum est atteint à un événement, et la première fois qu'il atteint ce maximum est une ouverture. Le critère de départage à égalité est là que se joue la règle semi-ouverte. Un ordre qui se ferme à 2 n'est plus actif à 2, donc il doit sortir avant qu'un ordre qui s'ouvre à 2 n'entre. Ici, cela ne demande rien de plus~: Python trie les tuples selon leur deuxième champ en cas d'égalité, et $-1$ se trie avant $+1$.

\iqwrong{Traiter d'abord les ouvertures en cas d'égalité des dates. L'exemple compte alors 2 à l'instant 2, alors que les deux ordres ne sont jamais vivants ensemble.}
\iqfollow{Et si les ordres arrivent déjà triés par date d'ouverture~?}{Garder un tas-min des dates de fermeture des ordres vivants. Avant chaque ouverture, retirer ceux dont la fermeture n'est pas postérieure à celle-ci, et la taille du tas donne le compte. Cela coûte encore $O(n \log n)$ dans le pire cas, mais la mémoire ne sert qu'aux ordres vivants à un instant donné, et la méthode fonctionne sur un flux.}
\iqfollow{Et si les dates sont des entiers dans un petit intervalle $0, \dots, R$~?}{Ajouter $+1$ à chaque ouverture et $-1$ à chaque fermeture dans un tableau de taille $R + 1$, puis prendre une somme cumulée~: $O(n + R)$, sans tri.}
\end{solution}

\begin{solution}{iq:iv3:algorithms:20}
Additionner le notionnel par symbole dans une table de hachage, puis prendre les $k$ premiers avec un tas de taille $k$, comme dans le \cref{lst:iv3:al:topk}. Le calcul des totaux coûte $O(n)$ en moyenne et la sélection $O(m \log k)$, avec une mémoire en $O(m)$.
\omcode{code/interviews-3/01-algorithms/python/iv_top_k.py}{1}{10}{Totaux, puis tas borné.}\label{lst:iv3:al:topk}

\texttt{heapq.nsmallest} avec la clé \texttt{(-notional, symbol)} conserve un tas des $k$ meilleurs candidats rencontrés jusque-là. Si vous écrivez le tas à la main, c'est un tas-min sur le notionnel, dont la racine est le candidat le plus faible conservé~: chaque nouveau symbole remplace la racine ou est écarté. L'ordre de départage doit alors être inversé dans ce tas, car la racine doit être la pire, c'est-à-dire ici le nom venant plus tard. Cette inversion est le bogue habituel. Pour $k$ proche de $m$, il est plus simple de trier, et \texttt{quickselect} trouve l'ensemble des $k$ premiers en $O(m)$ en moyenne, si l'ordre à l'intérieur de cet ensemble importe peu.

Pour les transactions \texttt{AAA 100 3}, \texttt{BBB 50 10}, \texttt{CCC 300 1}, \texttt{AAA 100 2} et \texttt{DDD 10 1}, les trois premiers sont AAA 500, BBB 500 et CCC 300.

Si les transactions continuent d'arriver, le tas ne suffit plus, car le total d'un symbole change alors qu'il se trouve dans le tas ou hors de celui-ci, et un tas ne permet ni de retrouver ni de relever une entrée. Une solution consiste à tenir un ensemble ordonné de couples (notionnel, symbole), en plus de la table~: chaque transaction retire puis réinsère une entrée, pour $O(\log m)$ par transaction, et la lecture des $k$ premiers coûte $O(k)$. Une autre consiste à empiler une nouvelle entrée à chaque mise à jour et à écarter les entrées périmées lorsqu'elles atteignent la racine.

\iqwrong{«~Un tas de taille $k$ fonctionne aussi sur le flux.~» Une fois qu'un symbole est sorti du tas, son total courant décide seul de son retour, donc il faut encore tous les totaux et une façon de comparer à la racine courante à chaque transaction.}
\end{solution}

\begin{solution}{iq:iv3:algorithms:21}
Garder un tas contenant une entrée par flux, qui porte la paire suivante de ce flux, comme dans le \cref{lst:iv3:al:merge}. Le coût est $O(N \log K)$ en temps et $O(K)$ en mémoire. Si un flux se tait, le générateur s'arrête dès que vient le tour de ce flux, et attend~: il ne peut pas savoir si un message plus ancien est encore en route.
\enlargethispage{-3\baselineskip}%
\omcode{code/interviews-3/01-algorithms/python/iv_merge_feeds.py}{1}{20}{Un tas contenant une entrée par flux.}\label{lst:iv3:al:merge}

La plus petite date restante se trouve toujours en tête de l'un des flux, donc dans le tas. Après avoir émis cette paire, on lit la paire suivante de ce flux et l'on remplace la racine, ce qui coûte $\log K$. Comme un flux n'a qu'une entrée au plus dans le tas, le couple (date, numéro de flux) est unique, et Python n'a jamais à comparer deux messages.

Le générateur ne lit que ce qui lui est nécessaire, il fonctionne donc sur des flux sans fin~: la fusion d'un flux qui émet toutes les 2 unités et d'un autre qui émet toutes les 3 unités donne 2, 3, 4, 6, 6, le flux de plus petit numéro passant en premier à égalité.

Le flux silencieux est la partie réelle du problème. Une fusion dans l'ordre temporel ne peut émettre un message que lorsque chaque flux a montré quelque chose de postérieur, si bien que le flux le plus lent fixe le rythme. La solution usuelle est un battement de cœur, que chaque flux envoie avec son temps courant pour que les autres puissent avancer, ou un délai d'attente au terme duquel on déclare le flux périmé et l'on continue sans lui, au prix d'un désordre possible à son retour.

\iqwrong{Empiler \texttt{(timestamp, message)}. Cela semble marcher jusqu'à ce que deux messages partagent une date. Python compare alors les dictionnaires et lève \texttt{TypeError}.}
\iqfollow{Concaténer puis trier~?}{$O(N \log N)$ et les $N$ éléments en mémoire, ce qui échoue sur un flux.}
\end{solution}

\begin{solution}{iq:iv3:algorithms:22}
Garder deux tas, comme dans le \cref{lst:iv3:al:median}~: un tas-max pour la moitié inférieure et un tas-min pour la moitié supérieure, dont les tailles diffèrent d'un élément au plus. Alors \texttt{add} coûte $O(\log n)$ et \texttt{median} coûte $O(1)$, avec une mémoire en $O(n)$. Pour des ticks entiers entre 0 et 1000, garder plutôt un tableau de 1001 compteurs.
\omcode{code/interviews-3/01-algorithms/python/iv_median_stream.py}{1}{17}{Deux tas.\protect\hfil\protect\break}\label{lst:iv3:al:median}

La médiane se trouve à la frontière entre les deux moitiés, seules les deux racines comptent. Chaque \texttt{add} passe d'abord par \texttt{high}, afin que le plus petit de la moitié supérieure bascule dans l'autre, puis redescend d'un cran si \texttt{low} est devenu trop grand. Ce seul parcours garantit les deux règles sans cas particulier~: chaque valeur de \texttt{low} est au plus égale à chaque valeur de \texttt{high}, et \texttt{high} porte l'élément en surplus lorsque l'effectif est impair. Le module \texttt{heapq} de Python ne fournit qu'un tas-min, donc \texttt{low} stocke des valeurs opposées, et le cas pair lit $(\text{high}_0 - \text{low}_0) / 2$. Pour les valeurs 5, 1, 9, 3, il donne 4.

Avec une plage bornée, un \texttt{add} est un simple incrément, en $O(1)$. La médiane peut alors se trouver en parcourant les compteurs, ou se suivre grâce à un pointeur qui se déplace à mesure que les valeurs arrivent~; la mémoire reste fixe, 1001 compteurs, quelle que soit la longueur du flux. Le gain est réel, car les tas gardent chaque valeur pour toujours.

Avant le premier \texttt{add} il n'existe pas de médiane, et le listing lève \texttt{IndexError}~: dites que ce choix vous appartient.

\iqwrong{Garder une liste triée avec une recherche dichotomique. La recherche coûte $O(\log n)$, mais l'insertion déplace jusqu'à $n$ éléments, si bien qu'un flux coûte $O(n^2)$. C'est acceptable pour un petit $n$, et il vaut la peine de le dire.}
\iqfollow{Et si les anciennes valeurs doivent aussi sortir~?}{Un tas ne sait pas retirer un élément du milieu, c'est donc un problème plus difficile~: une médiane glissante. Il faut une suppression paresseuse ou un arbre équilibré.}
\end{solution}

\begin{solution}{iq:iv3:algorithms:23}
Garder les $k$ premiers éléments. Ensuite, l'élément à la position $i$ (indexée à partir de 0) remplace un emplacement tiré uniformément avec probabilité $k/(i+1)$, comme dans le \cref{lst:iv3:al:reservoir}. C'est l'échantillonnage par réservoir\index{échantillonnage par réservoir}~: $O(n)$ en temps et $O(k)$ en mémoire, et chaque sous-ensemble de taille $k$ a la même probabilité d'être retenu.
\omcode{code/interviews-3/01-algorithms/python/iv_reservoir.py}{1}{10}{Échantillonnage par réservoir.}\label{lst:iv3:al:reservoir}

La preuve procède par récurrence sur le nombre $i$ d'éléments vus. Supposons qu'après $i$ éléments, chacun soit retenu avec probabilité $k/i$. L'élément $i+1$ entre avec probabilité $k/(i+1)$, comme le dit le code. Un élément antérieur survit à cette étape, sauf si le nouveau venu entre et choisit son emplacement, donc il est retenu avec probabilité
\[
\frac{k}{i}\Bigl(1 - \frac{k}{i+1} \cdot \frac{1}{k}\Bigr) = \frac{k}{i} \cdot \frac{i}{i+1} = \frac{k}{i+1}.
\]
Le cas de base est $i = k$, où tous sont retenus. Le test vérifie cette identité de façon exacte, et va plus loin~: pour $n = 5$ et $k = 2$, il énumère les 60 suites de tirages également probables, et chacun des 10 sous-ensembles apparaît exactement 6 fois. Une simulation à germe fixe confirme le résultat~: avec 20 éléments et $k = 5$, chacun est conservé environ un quart du temps. Les 20 valeurs concordent à environ un point de pourcentage près, sur 40\,000 essais.

\iqwrong{Remplacer avec probabilité $1/(i+1)$, la règle du cas $k = 1$. Le dernier élément est alors retenu avec probabilité $1/n$, alors qu'il faudrait $k/n$.}
\iqfollow{Et si les éléments ont des poids~?}{Donner à l'élément $i$ la clé $u_i^{1/w_i}$, où $u_i$ suit une loi uniforme sur 0 à 1, et garder les $k$ plus grandes clés dans un tas. Pour $k = 1$, l'élément $i$ survit avec probabilité $w_i / \sum_j w_j$~; la simulation confirme ce résultat, pour les poids 1, 2 et 5.}
\end{solution}

\begin{solution}{iq:iv3:algorithms:24}
Utiliser une liste doublement chaînée ordonnée par récence, et une table de hachage qui associe chaque clé à l'itérateur de la liste, comme dans le \cref{lst:iv3:al:lru}. \texttt{splice} rattache un nœud en tête en $O(1)$, sans toucher à aucun itérateur. Les deux opérations coûtent donc $O(1)$ en moyenne, et la mémoire est en $O(\text{capacité})$.
\omcode{code/interviews-3/01-algorithms/cpp/lru.hpp}{8}{39}{Une liste et une table de ses itérateurs.}\label{lst:iv3:al:lru}

Trois opérations doivent rester peu coûteuses~: trouver une clé, la placer en tête, et retirer la plus ancienne. La table de hachage assure la première, la liste les deux autres. La table stocke des itérateurs, ce qui n'est sûr que parce que \texttt{std::list} n'invalide jamais un itérateur sur un élément conservé, et \texttt{splice} non plus. Un \texttt{std::vector} décalerait les éléments à chaque déplacement, ce qui coûte $O(n)$. Un \texttt{std::deque} ne sait ajouter et retirer qu'à ses extrémités, donc déplacer un élément depuis le milieu coûte aussi $O(n)$, et l'insertion aux extrémités invalide tous ses itérateurs.

Le test lance des opérations aléatoires contre un modèle lent en Python, aussi sous AddressSanitizer et UndefinedBehaviorSanitizer, et les résultats concordent. L'éviction retire l'entrée de la table avant de dépiler le nœud, car la clé est lue dans ce nœud. Noter que \texttt{get} écrit~: il modifie la liste.

\iqwrong{«~Horodater chaque entrée et chercher la plus ancienne par un parcours.~» Les lectures sont rapides, mais l'éviction coûte $O(n)$. Un tas sur les horodatages échoue aussi, car chaque \texttt{get} devrait modifier une clé à l'intérieur du tas.}
\iqfollow{Est-ce sûr avec un thread si les lecteurs utilisent un verrou partagé~?}{Non, puisque \texttt{get} déplace un nœud, les lecteurs ont donc besoin d'un accès exclusif. Sous contention, découper le cache en fragments dotés chacun de leur verrou, ou mettre la récence à jour par lots.}
\iqfollow{Les nœuds de la liste sont dispersés en mémoire. Cela compte-t-il~?}{Chaque pas peut provoquer un défaut de cache, et chaque insertion une allocation. Une version pour le chemin critique (hot path) garde les nœuds dans un tableau de taille fixe et les relie par des indices entiers, avec la même logique.}
\end{solution}

\begin{solution}{iq:iv3:algorithms:25}
Lancer l'algorithme de Floyd à deux pointeurs\index{détection de cycle de Floyd}, comme dans le \cref{lst:iv3:al:cycle}~: l'un avance d'un pas à la fois, l'autre de deux. S'ils se rencontrent, il existe un cycle. Replacer alors l'un des deux en tête et faire avancer les deux d'un pas à la fois~: ils se retrouvent au début du cycle. Cela coûte $O(n)$ en temps et $O(1)$ en mémoire.
\omcode{code/interviews-3/01-algorithms/python/iv_cycle.py}{1}{10}{Deux vitesses, puis la même vitesse.}\label{lst:iv3:al:cycle}

Notons $a$ le nombre de nœuds de la queue avant le cycle, et $c$ la longueur du cycle. Supposons qu'ils se rencontrent après que le pointeur lent a parcouru $s$ pas. Le rapide en a parcouru $2s$, et les $s$ pas supplémentaires doivent former un nombre entier de tours, donc $s$ est un multiple de $c$. Le point de rencontre se trouve à $s$ pas de la tête. En avançant encore de $a$ pas à partir de là, on atteint la position $s + a$ depuis la tête, qui est le début du cycle, puis un nombre entier de tours. Un marcheur parti de la tête atteint le même nœud après $a$ pas, donc les deux se rencontrent là, et ils ne peuvent se rencontrer plus tôt, puisque le marcheur n'est pas sur le cycle avant ce moment.

\iqwrong{«~Stocker les nœuds visités dans un ensemble.~» C'est juste et le plus clair, mais cela coûte $O(n)$ en mémoire, ce que l'énoncé interdit.}
\iqfollow{Quelle est la longueur du cycle~?}{Une fois qu'ils se sont rencontrés, immobiliser l'un des deux et faire tourner l'autre jusqu'à son retour, en comptant.}
\end{solution}

\begin{solution}{iq:iv3:algorithms:26}
Trier par début, puis parcourir la liste en prolongeant la dernière pause fusionnée chaque fois que la suivante commence au plus tard à sa fin, comme dans le \cref{lst:iv3:al:intervals}. On obtient $(1, 6)$ et $(8, 12)$, soit 9 heures au total. Cela coûte $O(n \log n)$ en temps.
\omcode{code/interviews-3/01-algorithms/python/iv_intervals.py}{1}{12}{Trier par début, puis prolonger ou commencer.}\label{lst:iv3:al:intervals}

Après le tri, une nouvelle pause ne peut chevaucher aucune pause antérieure à la dernière fusionnée, car toutes les pauses précédentes se sont terminées avant le début de la dernière. Une seule comparaison suffit donc. Le \texttt{max} compte~: une pause courte incluse dans une longue, comme $(2, 3)$ après $(1, 6)$, ne doit pas réduire la fin. Le test \texttt{start <= end} décide qu'deux pauses qui se touchent ne forment qu'une seule pause~; le remplacer par \texttt{<} laisse $(8, 10)$ et $(10, 12)$ séparées.

Le total est la somme des longueurs fusionnées, 5 plus 4. Additionner les longueurs brutes donne 2 + 4 + 2 + 2 = 10, ce qui compte deux fois le recouvrement de 2 à 3.

\iqwrong{Trier par fin. Alors $(1, 10)$ vient après $(2, 3)$ et $(4, 5)$, et un simple retour sur la dernière pause fusionnée ne voit pas que la longue les absorbe toutes.}
\iqfollow{Et si les pauses arrivent une à une~?}{Garder les pauses fusionnées dans une table ordonnée, de début vers fin. Trouver les voisines par recherche dichotomique, et absorber celles qui se chevauchent~: $O(\log n)$ pour la recherche, et chaque pause n'est retirée qu'au plus une fois, donc le coût est en moyenne $O(\log n)$ par arrivée dans un arbre.}
\end{solution}

\begin{solution}{iq:iv3:algorithms:27}
Il renvoie \texttt{[1, 1, 4, 2, 1, 1, 0, 0]}. Le \cref{lst:iv3:al:nextgreater} garde une \index{pile monotone}pile monotone des jours qui attendent encore un prix plus élevé, et chaque nouveau prix règle chacun des jours en attente qu'il bat.
\omcode{code/interviews-3/01-algorithms/python/iv_next_greater.py}{1}{9}{Une pile de jours en attente.}\label{lst:iv3:al:nextgreater}

Les prix de la pile ne montent jamais du bas vers le haut, car tout prix qui bat le sommet le dépile avant de continuer. Si le nouveau prix ne bat pas le sommet, il ne bat rien de ce qui se trouve dessous, et la boucle peut s'arrêter.

Prenons 103 au jour 2~: il attend pendant que arrivent 99, 97 et 100, qui s'empilent ou se dépilent mutuellement, puis il est réglé au jour 6 par 104, ce qui donne 4.

Un seul prix peut dépiler plusieurs entrées, donc une étape peut coûter cher. Pourtant, chaque jour est empilé une fois et dépilé au plus une fois, si bien que sur toute la liste les dépilements sont au plus $n$~: $O(1)$ amorti par jour, $O(n)$ au total, et une mémoire en $O(n)$ dans le pire cas.

\iqwrong{Deux boucles imbriquées qui parcourent vers l'avant. Sur une liste décroissante, chaque jour parcourt jusqu'à la fin~: $n(n - 1)/2$ comparaisons.}
\iqfollow{Pour chaque jour, combien de jours consécutifs, se terminant aujourd'hui, ont eu un prix inférieur ou égal à celui d'aujourd'hui~?}{La même pile vue dans l'autre sens~: dépiler tout ce qui est inférieur ou égal au prix du jour, puis l'écart est la distance jusqu'au nouveau sommet. Pour \texttt{[100, 80, 60, 70, 60, 75, 85]}, on obtient 1, 1, 1, 2, 1, 4, 6.}
\end{solution}

\begin{solution}{iq:iv3:algorithms:28}
Le tri décroissant donne 4, 1, 1, soit trois ordres, alors que la vraie réponse est 3, 3, soit deux. Le \Cref{lst:iv3:al:clips} utilise la \index{programmation dynamique}programmation dynamique sur les quantités, en temps $O(ns)$ et mémoire $O(n)$ pour $s$ tailles.
\omcode{code/interviews-3/01-algorithms/python/iv_clips.py}{1}{15}{Nombre minimal d'ordres pour chaque quantité.}\label{lst:iv3:al:clips}

Un remplissage exact de $q$ se termine par une dernière taille $s$, et ce qui précède est un remplissage exact de $q - s$. Donc \texttt{best[q]} vaut un de plus que le plus petit \texttt{best[q - s]}. Si l'on retient la taille choisie pour chaque $q$, on reconstruit la liste en remontant. Si aucune combinaison ne fonctionne, \texttt{best[n]} garde sa valeur sentinelle, et avec les tailles 2 et 4 la fonction ne renvoie rien pour $n = 5$.

Notons que $ns$ n'est pas polynomial en la longueur de l'entrée, car $n$ s'écrit sur environ $\log n$ bits. Au tableau, la question est rarement assez grande pour que cela pose problème~; il faut toutefois le signaler si $n$ peut valoir $10^9$.

\iqwrong{«~Le tri décroissant est optimal, comme pour rendre la monnaie.~» C'est vrai pour certaines listes de tailles, comme 1, 5, 10 et 25, mais c'est faux pour celle-ci, et le tri décroissant peut même ne pas aboutir~: avec les tailles 3 et 4, 6 a pour réponse 3, 3, alors que le tri décroissant prend 4 et reste bloqué à 2.}
\iqfollow{De combien de façons peut-on faire cela, si l'ordre n'importe pas~?}{On boucle sur les tailles à l'extérieur et sur les quantités à l'intérieur, de sorte que chaque combinaison soit comptée une seule fois~: il y a 4 façons pour les tailles 1, 3, 4 et $n = 6$.}
\end{solution}

\begin{solution}{iq:iv3:algorithms:29}
La réponse est 9~: on achète à 10 et l'on vend à 18 pour un gain de 6, puis on achète à 14 et l'on vend à 19 pour un gain de 3. Le \Cref{lst:iv3:al:tradefee} conserve deux nombres, la meilleure richesse sans position et la meilleure richesse en position, mis à jour une fois par jour. Sans frais, elle renvoie la somme de toutes les hausses, ici 14~; avec des frais d'un million, elle renvoie 0.
\enlargethispage{-3\baselineskip}%
\omcode{code/interviews-3/01-algorithms/python/iv_trade_fee.py}{1}{6}{Deux états, un seul passage.}\label{lst:iv3:al:tradefee}

À la clôture de chaque jour, on est soit sans position, soit en détention d'une action, et la meilleure richesse de chaque état ne dépend que des deux valeurs de la veille. Pour détenir aujourd'hui, on détient déjà ou l'on achète depuis l'état sans position. Pour être sans position aujourd'hui, on l'était déjà ou l'on vend et l'on paie les frais. Aucune transaction n'est jamais forcée, donc \texttt{flat} part de 0 et ne peut que croître, ce qui explique qu'avec de gros frais on obtienne 0 et non une perte.

L'ordre des mises à jour est ici sûr~: utiliser le \texttt{holding} du jour dans le \texttt{flat} du jour n'autorise qu'un achat et une vente au même prix le même jour, ce qui coûte les frais et n'est donc jamais retenu.

\iqwrong{Négocier à chaque hausse et payer les frais à chaque fois. Ici cela donne $(3 - 2) + (6 - 2) + (5 - 2) = 8$, soit moins que 9, car garder la position pendant la baisse de 13 à 12 économise des frais et ne sacrifie qu'1 de prix.}
\iqfollow{Au plus deux allers-retours~?}{On conserve plutôt quatre états~: en position la première fois, sans position après la première, en position la seconde fois, sans position après la seconde.}
\end{solution}

\begin{solution}{iq:iv3:algorithms:30}
La distance de ABCD à ACBD est 2~: deux substitutions, car il n'existe pas ici d'opération d'échange. Le \Cref{lst:iv3:al:edit} remplit le tableau classique une ligne à la fois, en temps $O(mn)$ et mémoire $O(\min(m, n))$.
\omcode{code/interviews-3/01-algorithms/python/iv_edit.py}{1}{11}{Une seule ligne du tableau à la fois.}\label{lst:iv3:al:edit}

Soit $D(i, j)$ la distance entre les $i$ premiers caractères d'une chaîne et les $j$ premiers de l'autre. La dernière opération conserve ou substitue les derniers caractères, ce qui coûte 0 ou 1 en plus de $D(i - 1, j - 1)$, ou bien supprime ou insère un caractère, chacune au prix de 1 en plus de la case du dessus ou de celle de gauche. La première ligne et la première colonne sont les longueurs~: $D(i, 0) = i$.

Chaque case n'a besoin que de son voisin de gauche, de la case du dessus et de la diagonale. Une seule ligne, plus une variable sauvegardée, \texttt{diag}, contient les trois.

Le prix de cette économie est que l'on obtient le nombre et non les modifications. Pour retrouver les modifications à effectuer, on conserve le tableau complet, ou l'on utilise la méthode diviser pour régner de Hirschberg, qui demande une mémoire linéaire et environ deux fois plus de temps.

\iqwrong{Comparer des ensembles ou des lettres triées. ABCD et ACBD ont les mêmes lettres, donc ces méthodes les déclarent identiques.}
\iqfollow{Et si un échange de voisins coûte 1~?}{On ajoute un cas de plus, à partir de $D(i - 2, j - 2)$, lorsque les deux paires sont échangées. ABCD vers ACBD coûte alors 1. C'est pourquoi le rapprochement de noms utilise souvent cette variante.}
\end{solution}

\begin{solution}{iq:iv3:algorithms:31}
On utilise une \index{union-find}union-find, comme dans le \cref{lst:iv3:al:unionfind}~: chaque groupe est un arbre de pointeurs vers le parent, et deux comptes sont dans un même groupe s'ils atteignent la même racine. Avec une union par taille et une compression de chemin, $m$ opérations coûtent $O(m\,\alpha(n))$, où $\alpha$ est la fonction d'Ackermann inverse, qui est inférieure à 5 pour tout $n$ tenant en mémoire.
\omcode{code/interviews-3/01-algorithms/python/iv_union_find.py}{1}{24}{Pointeurs vers le parent, stockés à plat.}\label{lst:iv3:al:unionfind}

L'union par taille garde les arbres peu profonds~: la profondeur d'un compte n'augmente que lorsque son arbre est rattaché sous un arbre plus grand, si bien que son arbre au moins double à chaque fois, et la profondeur vaut au plus $\log_2 n$. La compression par moitié aplatit ensuite les nœuds traversés. Le nombre de groupes est un compteur qui diminue de un à chaque fusion effective.

Ce qu'elle ne peut pas faire, c'est scinder un groupe. Un lien ne peut pas être retiré, car la structure oublie quelle paire l'a créé. Si les suppressions sont connues à l'avance, on traite toute la séquence à rebours, ce qui les transforme en fusions.

\iqwrong{Étiqueter chaque compte avec le numéro de son groupe et réétiqueter à chaque fusion. Un test \texttt{same} coûte $O(1)$, et chaque fusion parcourt $O(n)$. Ne réétiqueter que le plus petit groupe donne $O(n \log n)$ pour toutes les fusions, ce qui suffit pour une réponse simple, bien que moins bon que les arbres.}
\iqfollow{Pourquoi ne pas faire un parcours en largeur à chaque requête~?}{Il répond correctement en $O(V + E)$ par requête. Si les requêtes sont fréquentes, c'est bien trop lent.}
\end{solution}

\begin{solution}{iq:iv3:algorithms:32}
\Cref{lst:iv3:al:topo} est le \index{tri topologique}tri topologique de Kahn~: on part des tâches qui n'attendent rien, et chaque fois qu'une tâche se termine, on libère celles qui n'attendaient plus qu'elle. Il faut $O(n + e)$ en temps pour $e$ paires. Si une tâche n'est jamais libérée, il existe un cycle. Le lot se termine à la 8\up{e} heure.
\omcode{code/interviews-3/01-algorithms/python/iv_topo.py}{1}{20}{Libérer les tâches dès que leur dernière dépendance se termine.\protect\hfil\protect\break}\label{lst:iv3:al:topo}

Un ordre existe exactement lorsque les paires ne forment pas de cycle. Chaque tâche d'un cycle attend une autre tâche de ce cycle, donc aucune d'elles n'atteint jamais un compteur d'attente égal à zéro. Réciproquement, si le graphe n'a pas de cycle, une tâche n'a toujours plus aucune dépendance non terminée, donc la boucle les termine toutes. Comparer la longueur de \texttt{order} à $n$ constitue donc à lui seul le test de cycle.

Avec un nombre illimité de travailleurs en parallèle, chaque tâche démarre dès que la dernière de ses dépendances se termine~: on parcourt donc les tâches dans l'ordre topologique et l'on prend le \texttt{max}. A se termine à 3. B se termine à 5 et C à 7. D démarre à 7 et se termine à 8. La chaîne A, C, D est le chemin critique, et accélérer B ne change rien.

\iqwrong{«~On additionne les durées.~» On obtient 10, qui est le temps sur un seul travailleur. Cela ne dit rien du cas parallèle.}
\iqfollow{Et si deux tâches seulement peuvent tourner à la fois~?}{L'ordonnancement avec un nombre limité de travailleurs est NP-difficile en général, on a donc recours à une heuristique, par exemple en lançant la tâche dont le chemin restant est le plus long.}
\end{solution}

\begin{solution}{iq:iv3:algorithms:33}
La réponse est 4 déplacements~: bas, bas, droite, droite. Le \Cref{lst:iv3:al:grid} utilise un \index{parcours en largeur}parcours en largeur, en temps et en mémoire $O(rc)$.
\enlargethispage{-3\baselineskip}%
\omcode{code/interviews-3/01-algorithms/python/iv_grid.py}{1}{17}{Parcours en largeur avec marquage à l'empilement.}\label{lst:iv3:al:grid}

La file contient les cases dans l'ordre de leur distance~: tout ce qui est à la distance $d$ sort avant ce qui est à $d + 1$. La première fois qu'une case est atteinte, c'est donc par un plus court chemin, et elle n'a plus besoin d'être revue. Le parcours en profondeur n'a pas un tel ordre. Sur une grille libre, il peut s'engager d'abord dans un long chemin et marquer des cases comme visitées~; le premier chemin qu'il trouve ne dit alors rien du plus court, et le corriger reviendrait à essayer tous les chemins.

Le marquage compte en revanche pour le coût, non pour la justesse. Marquée au dépilement, une case peut être empilée une fois par chacun de ses quatre voisins au plus avant d'être dépilée pour la première fois. Le résultat reste correct, mais la file est plus grande, jusqu'à quatre fois, et chaque répétition reparcourt ses voisins.

\iqwrong{«~On utilise Dijkstra.~» L'algorithme fonctionne, mais il ajoute un facteur $\log$ pour rien, puisque chaque déplacement coûte la même chose.}
\iqfollow{Certains déplacements coûtent 0 et d'autres 1~?}{On utilise une file à deux extrémités~: on empile en tête une case atteinte au coût 0 et en queue une case atteinte au coût 1. On conserve le même ordre, sans tas, et le temps reste linéaire.}
\end{solution}

\begin{solution}{iq:iv3:algorithms:34}
Le chemin le plus rapide de A à D prend 4 microsecondes, par A, C, B, D. Le \Cref{lst:iv3:al:dijkstra} est l'\index{algorithme de Dijkstra}algorithme de Dijkstra avec un tas binaire, en temps $O((V + E) \log V)$. Les prédécesseurs donnent le chemin, que l'on remonte depuis la cible.
\omcode{code/interviews-3/01-algorithms/python/iv_dijkstra.py}{1}{16}{Dijkstra avec suppression paresseuse.}\label{lst:iv3:al:dijkstra}

Le tas de Python ne permet pas d'abaisser une clé~: chaque amélioration empile donc une nouvelle entrée, et l'entrée laissée derrière est ignorée lorsqu'elle remonte. C'est le rôle de la ligne \texttt{d > dist[v]}. Le tas contient alors jusqu'à $E$ entrées, ce qui ne change rien à la borne, puisque $\log E$ est au plus $2 \log V$.

La preuve repose sur des poids non négatifs~: lorsque la plus petite entrée sort, aucune entrée encore en attente ne peut l'atteindre plus économiquement, car tout autre chemin passe par une entrée au moins aussi grande. Un arc négatif rompt cet argument. Avec A vers B de coût 2, A vers C de coût 3 et C vers B de coût $-2$, la version à ensemble fermé fixe B à 2, alors que la vraie réponse est 1. La version paresseuse ci-dessus trouve bien 1 lorsqu'un nœud est revisité, mais son temps d'exécution peut alors exploser, et si un cycle de coût négatif est accessible, la boucle ne se termine jamais. Pour les arcs négatifs, on utilise Bellman--Ford, en $O(VE)$.

\iqwrong{«~Parcours en largeur.~» Il trouve le plus petit nombre d'arcs. Ici, cela donne A, B, D à 6 microsecondes, avec un saut de moins mais un temps moins bon.}
\iqfollow{Une seule cible~?}{On s'arrête lorsqu'elle sort du tas. Sa distance est définitive à cet instant.}
\end{solution}

\begin{solution}{iq:iv3:algorithms:35}
La réponse est 16 pour $d = 2$, avec le découpage 4, 9, 3 et 7, 8, et 13 pour $d = 3$, avec 4, 9 et 3, 7 et 8. Le \Cref{lst:iv3:al:splitdays} utilise une \index{recherche dichotomique sur la réponse}recherche dichotomique sur la réponse, en temps $O(n \log S)$, où $S$ est la somme.
\enlargethispage{-3\baselineskip}%
\omcode{code/interviews-3/01-algorithms/python/iv_split_days.py}{1}{17}{Recherche dichotomique sur la plus grosse journée.}\label{lst:iv3:al:splitdays}

On retourne la question. Étant donné un plafond par jour, peut-on tout faire en $d$ jours~? La réponse est simple~: on remplit chaque jour de manière gloutonne jusqu'à ce que la tranche suivante dépasse le plafond. Prendre le plus possible aujourd'hui ne nuit jamais, puisqu'il en reste moins pour plus tard. Et la réponse est monotone~: si un plafond convient, tout plafond plus grand convient aussi. C'est ce qui permet de faire une recherche dichotomique sur le plafond, entre la plus grande tranche et la somme.

Des tailles positives rendent le test glouton correct. Avec une tranche négative, en ajouter une de plus peut faire baisser le total d'une journée, si bien que le remplissage glouton n'est plus optimal, et son nombre de jours peut même augmenter avec le plafond~: sur \texttt{[8, 2, -5, 8, -1, 4]}, les plafonds 8, 9, 10 et 11 donnent 2, 2, 3 et 2 jours. La recherche dichotomique n'a alors plus rien de solide à tester. La plus grande tranche ne constitue plus non plus une borne inférieure, car une journée peut contenir moins que l'une de ses tranches.

\iqwrong{«~On vise la moyenne, la somme divisée par $d$, et l'on coupe dès qu'on la dépasse.~» Pour $d = 2$, la moyenne vaut 15.5~; couper avant de la dépasser donne 13 puis 18, si bien que la plus grande journée vaut 18, contre 16 pour la vraie réponse.}
\iqfollow{L'alternative classique~?}{La programmation dynamique sur (position, jours utilisés), en $O(n^2 d)$, qui donne aussi les coupes, et qui est nécessaire si le coût d'une journée n'est pas la somme.}
\end{solution}

\begin{solution}{iq:iv3:algorithms:36}
\index{quickselect}Quickselect~: on partitionne autour d'un pivot aléatoire, et l'on ne garde que le côté qui contient la position $k$, comme dans le \cref{lst:iv3:al:quickselect}. En moyenne, cela coûte $O(n)$, car le travail décroît à peu près géométriquement. Trois groupes (inférieurs, égaux, supérieurs) et un pivot aléatoire évitent les deux mauvais cas.
\omcode{code/interviews-3/01-algorithms/python/iv_quickselect.py}{1}{21}{Trois groupes autour d'un pivot aléatoire.}\label{lst:iv3:al:quickselect}

Avec le premier élément comme pivot sur une liste triée, chaque étape ne retire qu'un seul élément. Trouver le plus grand de 100 coûte alors $99 + 98 + \dots + 1 = 4950$ comparaisons, soit $n(n-1)/2$.

Si toutes les valeurs sont égales et que l'on n'a que deux groupes, chaque valeur tombe dans «~pas inférieure au pivot~». Là encore, seul le pivot sort à chaque étape, et 100 valeurs égales à la recherche de la position 99 coûtent aussi 4950 comparaisons. La partition à trois voies s'arrête tout de suite, puisque toute la liste forme le groupe «~égal~» et que $k$ tombe dans celui-ci, après un seul passage.

Un pivot aléatoire rend le mauvais cas peu probable, quelle que soit l'entrée, si bien qu'aucun adversaire ni aucune liste triée ne peut le forcer. Si l'on tient absolument à une garantie dans le pire cas, la médiane des médianes donne $O(n)$ dans le pire cas, avec une constante plus grande.

\iqwrong{«~Utiliser un tas de taille $k$.~» C'est en $O(n \log k)$, ce qui convient et fonctionne sur un flux, mais cela ne respecte pas la borne en temps moyen demandée ici.}
\iqfollow{Est-ce en place~?}{Le listing copie par souci de clarté. Une version réelle partitionne en place, et c'est ce que fait \texttt{std::nth\_element}. Elle laisse la liste réordonnée, ce qu'un appelant ne souhaite pas forcément.}
\end{solution}

\begin{solution}{iq:iv3:algorithms:37}
Elle renvoie $(-2, 0, 2)$, $(-2, 1, 1)$, $(-1, -1, 2)$ et $(-1, 0, 1)$. Le \Cref{lst:iv3:al:threesum} trie, fixe la plus petite valeur du triplet, puis parcourt le reste avec deux pointeurs qui convergent vers l'intérieur. Le temps est en $O(n^2)$ et la mémoire supplémentaire, hors tri et sortie, en $O(1)$.
\omcode{code/interviews-3/01-algorithms/python/iv_three_sum.py}{1}{19}{Trier, fixer une valeur, puis deux pointeurs.}\label{lst:iv3:al:threesum}

Pour une première valeur fixée, la paire doit avoir pour somme son opposé. Sur un intervalle trié, une somme trop petite ne peut être corrigée qu'en remontant le pointeur gauche, et une somme trop grande en descendant le pointeur droit, de sorte que chaque pointeur bouge au plus $n$ fois.

Les répétitions s'introduisent à deux endroits, et le code les arrête tous les deux. Des premières valeurs égales donnent les mêmes triplets, si bien que toutes sauf la première sont ignorées. Et après une réussite, des secondes valeurs égales redonnent la même paire, donc \texttt{lo} les dépasse. Le pointeur droit n'a alors besoin d'aucune précaution particulière~: dès que \texttt{lo} a changé de valeur, le même \texttt{hi} donne une somme différente et l'on avance.

\iqwrong{Trois boucles imbriquées avec un ensemble de triplets triés. C'est correct, en $O(n^3)$. Un ensemble masque aussi la question de l'origine des répétitions.}
\iqfollow{Avec un ensemble de hachage à la place des deux pointeurs~?}{Également en $O(n^2)$, cela fonctionne sans tri, mais garder les triplets distincts est plus délicat, et cela coûte $O(n)$ de mémoire supplémentaire.}
\end{solution}

\begin{solution}{iq:iv3:algorithms:38}
La réponse est 5, pour la suite C, C, D, D, D. Le \Cref{lst:iv3:al:kdistinct} fait grandir une \index{fenêtre glissante}fenêtre glissante vers la droite et la fait rétrécir par la gauche dès qu'elle contient trop de symboles, en $O(n)$ de temps et $O(k)$ de mémoire.
\omcode{code/interviews-3/01-algorithms/python/iv_k_distinct.py}{1}{15}{Une fenêtre avec un compteur par symbole.}\label{lst:iv3:al:kdistinct}

Une suite comportant au plus $k$ symboles distincts ne contient que de telles suites, si bien que la règle est monotone~: pour chaque extrémité droite, il existe un début gauche le plus à gauche valide, et celui-ci ne recule jamais lorsque l'extrémité droite avance. C'est pourquoi les pointeurs n'avancent que vers l'avant, et chacun bouge au plus $n$ fois. L'entrée d'un symbole est supprimée dès que son compteur atteint 0, puisque \texttt{len(count)} est ce que l'on compare à $k$.

La fenêtre est une suite valide à chaque étape, donc la plus longue rencontrée est la réponse, et elle fournit aussi le début gauche le plus à gauche pour chaque fin~: tous les débuts ultérieurs jusqu'à la fin sont valides eux aussi. C'est ce qui répond à la relance.

\iqwrong{Compter exactement $k$ avec une seule fenêtre directement. Lorsque la fenêtre contient exactement $k$ symboles, on ne voit pas quels débuts comptent, car plusieurs bords gauches peuvent aussi donner exactement $k$.}
\iqfollow{Et compter les suites comportant exactement $k$ symboles distincts~?}{À chaque extrémité droite, le nombre de suites valides comportant au plus $k$ symboles vaut \texttt{right - left + 1}. Les additionner, puis retrancher la somme pour au plus $k - 1$ symboles de celle pour au plus $k$. Pour la bande ci-dessus et $k = 2$, cela donne 11.}
\end{solution}

\begin{solution}{iq:iv3:algorithms:39}
Le \Cref{lst:iv3:al:flatmap} est une telle table, avec un \index{adressage ouvert}adressage ouvert. Vider une case lors d'un \texttt{erase} casse \texttt{get}, car une clé ultérieure qui a sondé au-delà de cette case semble maintenant absente. La correction retenue ici est la \index{suppression par décalage arrière}suppression par décalage arrière~: on ramène les entrées suivantes dans le trou, lorsque cela est permis.
\omcode{code/interviews-3/01-algorithms/cpp/flat_map.hpp}{6}{45}{Table plate avec décalage arrière.}\label{lst:iv3:al:flatmap}

\texttt{get} s'arrête à la première case vide, ce qui n'est sûr que si aucune clé ne se trouve au-delà d'un trou dans sa séquence de sondage. Après un simple vidage, cette règle n'est plus respectée. Les solutions usuelles sont une pierre tombale, c'est-à-dire un marqueur que \texttt{get} franchit et que \texttt{put} peut réutiliser, mais qui impose un rehachage périodique car les pierres tombales s'accumulent, ou le décalage utilisé ici. Une clé située à la case \texttt{j} ne peut revenir dans le trou en \texttt{i} que si sa case d'origine ne se trouve pas, cycliquement, entre les deux, ce que vérifie le test de distance. Avec un type non signé et un masque, le débordement circulaire ne demande rien de plus.

Le fait d'être plein est une question de longueurs de sondage. Près du facteur de charge $\alpha$, une recherche infructueuse avec sondage linéaire inspecte environ $\frac{1}{2}(1 + 1/(1 - \alpha)^2)$ cases. Dans une simulation à $2^{16}$ cases, cela fait environ 2.5 à une charge de 0.5, et environ 50 à une charge de 0.9. À charge 1, une recherche d'une clé absente ne se termine jamais ici, car il n'existe aucune case vide où s'arrêter. Il faut donc maintenir la charge autour de la moitié.

\iqwrong{«~Utiliser \texttt{std::unordered\_map}.~» C'est correct, mais cela alloue un nœud par clé, et les défauts de cache sont précisément ce que cette conception évite.}
\end{solution}

\begin{solution}{iq:iv3:algorithms:40}
Le \Cref{lst:iv3:al:heap} range l'arbre dans un tableau, les enfants de $i$ se trouvant en $2i + 1$ et $2i + 2$. \texttt{push} et \texttt{pop} coûtent $O(\log n)$, le minimum est à la racine en $O(1)$, et le constructeur coûte $O(n)$.
\omcode{code/interviews-3/01-algorithms/python/iv_heap.py}{1}{28}{Un tas dans un tableau.}\label{lst:iv3:al:heap}

\texttt{push} place l'élément en fin de tableau et le fait remonter par échanges. \texttt{pop} déplace le dernier élément à la racine et le fait descendre par échanges avec le plus petit des enfants. Comparer avec le plus petit enfant est essentiel~: avec le plus grand, cet enfant pourrait se retrouver au-dessus de l'autre et l'ordre serait rompu.

La construction est linéaire à cause de la répartition du travail. La moitié des nœuds sont des feuilles et ne descendent rien. Un nœud de hauteur $h$ descend au plus de $h$ niveaux, et il existe au plus $\lceil n/2^{h+1} \rceil$ nœuds de hauteur $h$, si bien que le total vaut environ $\sum_h h\,n/2^{h+1} = n$. Remonter par insertions successives est l'inverse~: la plupart des nœuds se trouvent près du bas, et chacun peut remonter jusqu'en haut. Sur une liste décroissante de 1024 éléments, les insertions font 8\,204 échanges, alors que la construction en fait 1\,015.

\iqwrong{«~La construction fait $n/2$ appels de coût $\log n$, donc elle est en $O(n \log n)$.~» C'est une borne supérieure vraie, mais pas serrée. Elle ignore que la plupart des appels se trouvent près du bas et sont peu coûteux.}
\iqfollow{Comment abaisser une clé déjà présente dans le tas~?}{Tenir une table qui associe chaque élément à sa position, mise à jour à chaque échange, puis faire remonter l'élément. Un tas dans un tableau n'a pas d'autre moyen de retrouver un élément.}
\end{solution}

\begin{solution}{iq:iv3:algorithms:41}
Elle renvoie 9 et 2. Le \Cref{lst:iv3:al:twosingles} fait le OU exclusif de tous les éléments pour obtenir $a \oplus b$, choisit un bit sur lequel ils diffèrent, puis refait le OU exclusif sur les seuls identifiants ayant ce bit à 1. Le temps est en $O(n)$ et la mémoire en $O(1)$.
\omcode{code/interviews-3/01-algorithms/python/iv_two_singles.py}{1}{10}{OU exclusif, découpage sur un bit, puis nouveau OU exclusif.}\label{lst:iv3:al:twosingles}

Le OU exclusif annule les paires, si bien que le premier passage laisse $9 \oplus 2$, soit en binaire $1001 \oplus 0010 = 1011$, qui vaut 11. Ce résultat n'est pas nul parce que les deux nombres diffèrent. Un bit à 1 de ce résultat, ici le plus bas, \texttt{both \& -both}, qui vaut 1, partage les identifiants en deux groupes. Chaque identifiant qui apparaît deux fois tombe dans un groupe avec son jumeau et s'annule de nouveau, et les deux singletons tombent dans des groupes différents, par le choix du bit. Le OU exclusif d'un groupe est donc l'un des singletons, et $a \oplus b \oplus a$ donne l'autre.

L'astuce \texttt{x \& -x} fonctionne parce que la négation en complément à deux inverse tous les bits situés au-dessus du bit à 1 le plus bas, et conserve ce bit.

\iqwrong{Utiliser ensemble la somme et le OU exclusif. Ces deux informations ne déterminent pas la paire~: 1 et 6, 2 et 5, 3 et 4 ont tous pour somme 7 et pour OU exclusif 7. Le découpage sur un bit est ce qui fournit une preuve.}
\iqfollow{Chaque identifiant apparaît trois fois, sauf un~?}{Le OU exclusif ne peut pas annuler trois copies. Compter chaque position de bit modulo 3~: un reste de 1 signifie que le singleton possède ce bit. Cela coûte $O(32n)$ de temps avec une mémoire constante, et donne 5 pour \texttt{[7, 7, 7, 5, 9, 9, 9]}.}
\end{solution}

\begin{solution}{iq:iv3:algorithms:42}
Les réponses sont oui, oui, oui, non, oui. Le \Cref{lst:iv3:al:ratelimit} conserve les instants des messages acceptés encore présents dans la fenêtre. Cela coûte une mémoire en $O(N)$ et un temps amorti en $O(1)$, et la limite n'est jamais dépassée.
\enlargethispage{-3\baselineskip}%
\omcode{code/interviews-3/01-algorithms/python/iv_rate_limit.py}{1}{17}{Une fenêtre glissante exacte.}\label{lst:iv3:al:ratelimit}

À l'instant 10, le message émis à 0 est sorti de la fenêtre, si bien que 1 et 10 l'occupent, et le second message à 10 est refusé. À 11, le message émis à 1 est sorti, il y a donc de nouveau de la place. Les messages refusés ne sont pas enregistrés, car ils n'ont pas atteint la plateforme.

Un compteur remis à zéro toutes les $T$ secondes semble convenir, et pourtant il peut doubler le débit à une frontière. Avec une remise à zéro à l'instant 10, les messages à 9, 9, 10 et 10 passent tous, soit 4 dans une fenêtre de 2 secondes, pour une limite de 2. Le prix de la violation de la limite est fixé par la règle propre à la plateforme~: il faut donc lire sa spécification.

Un \index{seau à jetons}seau à jetons ne conserve qu'un solde et un instant, il utilise donc une mémoire en $O(1)$, ce qui convient à un chemin critique. Mais le seau autorise une rafale complète de $N$ jetons plus le réapprovisionnement, si bien qu'à un débit de $N/T$ il peut envoyer jusqu'à environ $2N$ dans une fenêtre de $T$, à moins de le dimensionner pour rester sous la règle de la plateforme~: prendre $N/2$ de chaque est la manière usuelle de le rendre sûr.

\iqwrong{«~Retarder plutôt que refuser.~» C'est un choix distinct, fait après ce contrôle. Une file d'ordres retardés doit aussi pouvoir être annulée, afin de ne pas envoyer un prix devenu périmé.}
\end{solution}

\begin{solution}{iq:iv3:algorithms:43}
Conserver, pour chaque clé, un tableau croissant d'instants et un tableau de valeurs. \texttt{set} ajoute en fin de tableau, en temps amorti $O(1)$, et \texttt{get} est une recherche dichotomique en $O(\log m)$ pour $m$ entrées sous cette clé, comme dans le \cref{lst:iv3:al:asof}. La mémoire est proportionnelle à l'ensemble des écritures.
\omcode{code/interviews-3/01-algorithms/python/iv_as_of.py}{1}{16}{Recherche dichotomique sur les instants de chaque clé.}\label{lst:iv3:al:asof}

\texttt{bisect\_right} renvoie la position qui suit le dernier instant inférieur ou égal à $t$, si bien que l'entrée juste avant est la réponse, et une position de 0 signifie qu'il n'en existe aucune. Une valeur fixée exactement à $t$ compte comme «~inférieure ou égale~», et choisir \texttt{bisect\_left} à cet endroit est l'erreur habituelle~: elle renvoie la valeur précédente. Par exemple, avec des valeurs aux instants 1 et 3, \texttt{get} à 3 donne la seconde et \texttt{get} à 2 donne la première.

Des instants arrivant dans le désordre rompent l'ajout bon marché en fin de tableau. \texttt{bisect.insort} conserve l'ordre, et chaque insertion coûte $O(m)$ pour le décalage, ce qui convient si les arrivées tardives sont rares et tombent près de la fin. Si elles sont fréquentes, utiliser un dictionnaire ordonné, en $O(\log m)$ pour les deux opérations. Dans tous les cas, préciser ce qu'une insertion change à une réponse déjà donnée~: une requête pour un instant passé peut désormais renvoyer une autre valeur, et le code de recherche qui a besoin de rejouer les événements doit aussi stocker le moment où chaque valeur a été connue.

\iqwrong{«~Parcourir depuis la fin.~» C'est correct pour une requête à l'instant le plus récent, mais c'est linéaire pour n'importe quel instant ancien.}
\iqfollow{Joindre deux séries triées à chaque transaction~?}{Parcourir les deux avec deux pointeurs, en retenant la dernière cotation vue, en $O(m + n)$~: aucune recherche dichotomique n'est nécessaire. Des transactions aux instants 1, 2 et 5, avec des cotations aux instants 2 et 4, n'obtiennent aucune cotation pour la première, la cotation de 2 pour la deuxième, et la cotation de 4 pour la troisième.}
\end{solution}

\begin{solution}{iq:iv3:algorithms:44}
Conserver une liste des éléments, plus un dictionnaire qui associe chaque élément à sa position, comme dans le \cref{lst:iv3:al:randomset}. \texttt{pick} tire une position uniforme dans la liste. Les trois opérations sont en $O(1)$ en moyenne.
\omcode{code/interviews-3/01-algorithms/python/iv_random_set.py}{1}{25}{Une liste et un dictionnaire de positions.}\label{lst:iv3:al:randomset}

Un ensemble de hachage seul n'a aucun moyen de désigner son $i$-ème élément, si bien qu'un tirage uniforme exige un parcours, ou alors une marche aléatoire sur les compartiments, qui favorise les éléments seuls dans leur compartiment. Une liste seule tire bien, mais elle exige un parcours pour retirer un élément.

L'astuce pour le retrait est que la liste n'a pas besoin de rester ordonnée. Déplacer le dernier élément dans le trou coûte une étape et ne laisse aucun vide.

L'erreur classique se trouve en fin de liste. Si $x$ est le dernier élément, une version naïve retire le dernier et le réécrit à la position $i$, qui se trouve alors au-delà de la fin, ce qui provoque soit une erreur, soit une entrée périmée. Ici, \texttt{last != x} protège ce cas. L'autre erreur classique consiste à oublier de mettre à jour, dans le dictionnaire, la position de l'élément déplacé.

L'uniformité se vérifie par simulation dans le test~: après une suite aléatoire d'insertions et de retraits, chacun des éléments encore présents dans l'ensemble sort avec à peu près la même fréquence.

\iqfollow{Et si chaque élément a un poids~?}{La liste ne sert plus à rien. Avec des poids fixes, utiliser la méthode d'alias. S'ils changent, un arbre de sommes partielles donne $O(\log n)$ par opération.}
\end{solution}

\begin{solution}{iq:iv3:algorithms:45}
Le \Cref{lst:iv3:al:shuffle} est le \index{Fisher--Yates, mélange de}mélange de Fisher--Yates, qui tire chaque position uniquement parmi la partie qui n'est pas encore fixée. La version de votre collègue a $3^3 = 27$ issues équiprobables, et 27 n'est pas divisible par les $3! = 6$ ordres, si bien qu'elle ne peut pas être équitable.
\enlargethispage{-3\baselineskip}%
\omcode{code/interviews-3/01-algorithms/python/iv_shuffle.py}{1}{7}{Fisher--Yates, depuis la fin.}\label{lst:iv3:al:shuffle}

Fisher--Yates effectue $n$ choix, puis $n - 1$, et ainsi de suite jusqu'à 2~: cela fait $n!$ chemins équiprobables au total, un pour chaque ordre, si bien que chaque ordre sort avec la probabilité $1/n!$. La dernière position reçoit chaque élément avec la probabilité $1/n$, et il en va de même pour ce qui reste.

Pour la version de votre collègue, on compte par énumération. En partant de \texttt{1, 2, 3}, l'ordre 1, 3, 2 sort 5 fois sur 27, de même que 2, 1, 3 et 2, 3, 1. Les ordres 1, 2, 3~; 3, 1, 2 et 3, 2, 1 sortent chacun 4 fois. Le biais est donc de $5/27$ contre $4/27$. Il ne s'atténue pas pour un $n$ plus grand, car $n^n$ n'est jamais divisible par $n!$ pour $n \ge 3$~: $n - 1$ divise $n!$ et n'a aucun facteur premier commun avec $n^n$.

\iqwrong{«~Attribuer une clé aléatoire à chaque élément, puis trier.~» C'est équitable si les clés ne sont jamais égales, mais le coût est en $O(n \log n)$. Un comparateur qui renvoie une réponse aléatoire n'est pas équitable, et il peut casser le tri.}
\iqfollow{Un seul générateur aléatoire suffit-il pour un $n$ grand~?}{Un état sur 64 bits n'a que $2^{64}$ suites possibles, et $21!$ dépasse déjà ce nombre, si bien que la plupart des ordres, même d'un jeu de cartes court, ne peuvent pas sortir.}
\end{solution}

\begin{solution}{iq:iv3:algorithms:46}
On parcourt l'arbre en portant l'intervalle ouvert dans lequel chaque clé doit se trouver, comme dans le \cref{lst:iv3:al:bst}. Aller à gauche fixe une nouvelle borne supérieure, et aller à droite une nouvelle borne inférieure. Le temps est en $O(n)$, et la mémoire en $O(h)$ pour un arbre de hauteur $h$. Une pile à la place de la récursion gère 100\,000 niveaux.
\omcode{code/interviews-3/01-algorithms/python/iv_bst.py}{1}{13}{Chaque nœud dans un intervalle hérité de ses ancêtres.}\label{lst:iv3:al:bst}

La condition porte sur des sous-arbres entiers, et les enfants directs d'un nœud ne la couvrent pas. Prenons la racine 5, avec pour enfant gauche 1, dont l'enfant droit est 6, et 6 a pour enfant gauche 4. Chaque nœud est correct par rapport à ses enfants, et pourtant 4 se trouve dans le sous-arbre droit de 5, si bien que l'arbre est invalide. Dans le listing, au moment d'atteindre 4, son intervalle va de 5 à 6, et 4 tombe en dehors.

L'autre voie classique est un parcours infixe, qui doit visiter les clés dans l'ordre strictement croissant. Il n'a besoin que de la clé précédente, il est correct, et il demande le même temps et la même mémoire.

En Python, la récursion s'arrête par défaut vers une profondeur de 1000, et en C++ un arbre aussi profond peut faire déborder la pile, selon sa taille. C'est pourquoi le listing gère sa propre pile.

\iqwrong{«~Comparer avec les enfants~», ce qui accepte l'arbre ci-dessus. C'est la réponse fausse la plus fréquente, car elle paraît juste sur chaque arbre valide tracé au tableau.}
\iqfollow{La clé de rang $k$ dans l'ordre croissant~?}{On parcourt dans l'ordre et l'on s'arrête après $k$ clés~: $O(h + k)$. Si la question revient souvent, on garde la taille de chaque sous-arbre dans son nœud, pour un coût en $O(h)$.}
\end{solution}

\begin{solution}{iq:iv3:algorithms:47}
Le \Cref{lst:iv3:al:welford} est la \index{Welford, méthode de}méthode de Welford~: elle met à jour la moyenne et la somme des carrés des écarts à celle-ci, en $O(1)$ par valeur. Sur les quatre valeurs, elle donne 30. La méthode fondée sur le compte, la somme et la somme des carrés donne $-170.67$ en double précision, une variance négative.
\omcode{code/interviews-3/01-algorithms/python/iv_welford.py}{1}{12}{Mise à jour de Welford.}\label{lst:iv3:al:welford}

La formule $(\sum x^2 - (\sum x)^2/n) / (n - 1)$ soustrait deux nombres proches de $4 \times 10^{18}$, qui est bien supérieur à $2^{53}$, si bien que les doubles y sont espacés d'environ 512. La différence est voisine de 90, donc ce qui reste est surtout de l'erreur d'arrondi. La méthode de Welford travaille sur les écarts à la moyenne courante, qui sont petits, de sorte que peu d'information est perdue. Le produit $d \times (x - \text{moyenne})$ utilise la moyenne avant et après la mise à jour, il n'est jamais négatif, et la somme reste donc non négative.

Un cas limite~: si toutes les valeurs sont égales, chaque écart vaut 0, si bien que la variance vaut exactement 0, quelle que soit leur taille.

\iqwrong{«~Soustraire la première valeur de toutes les autres.~» Cela aide sur cet exemple, mais il faut conserver cette valeur, et un flux qui s'en éloigne fait revenir le problème.}
\iqfollow{Et avec une pondération exponentielle~?}{Avec un facteur d'oubli $\lambda$, on pose $d = x - \text{moyenne}$. Alors $\text{moyenne} \mathrel{+}= (1 - \lambda)\,d$ et $\text{var} = \lambda(\text{var} + (1 - \lambda)\,d^2)$. Chaque mise à jour est en $O(1)$, et aucune fenêtre n'est conservée.}
\iqfollow{Et une fenêtre des $k$ dernières valeurs~?}{On garde aussi les $k$ valeurs, et l'on retire la plus ancienne par la mise à jour inverse. Les arrondis s'accumulent, donc il faut recalculer de temps en temps.}
\end{solution}

\begin{solution}{iq:iv3:algorithms:48}
Le \Cref{lst:iv3:al:varint} encode les différences 1000, 1, $-2$ et 301. Par le codage en zigzag, elles deviennent 2000, 2, 3 et 602, qui occupent 2, 1, 1 et 2 octets~: 6 octets au total, contre 32 pour quatre entiers 64 bits non codés.
\enlargethispage{-3\baselineskip}%
\omcode{code/interviews-3/01-algorithms/python/iv_varint.py}{1}{23}{Différences, zigzag et varint.}\label{lst:iv3:al:varint}

Un varint consacre un octet aux valeurs inférieures à 128, puis un octet de plus à chaque tranche de sept bits supplémentaire. Ainsi 2 et 3 tiennent chacun sur un octet, tandis que 602 en demande deux, car il est supérieur ou égal à 128 et inférieur à 16\,384.

Le zigzag existe pour les petites valeurs négatives. Écrit comme un entier en complément à deux sur 64 bits, $-2$ a tous ses bits de poids fort à 1, si bien que, en varint non signé, il occupe dix octets, le plus long de tous. Le zigzag replie la droite numérique, de sorte qu'une petite valeur absolue donne un petit code quel que soit le signe, et $-2$ tient sur un seul octet.

Le décodeur inverse chaque étape~: il lit sept bits à la fois jusqu'à un octet dont le bit de poids fort est à 0, défait le zigzag, puis ajoute la valeur au total courant. Si un prix ne change pas, la différence vaut 0, un seul octet à 0. La méthode est donc à son meilleur sur les prix calmes.

\iqwrong{«~Compresser plutôt avec un compresseur générique.~» Cela aide, et les deux approches sont complémentaires, car les différences rendent les données plus compressibles. Mais cela coûte bien plus cher par message, et ne permet pas de décoder une seule valeur à faible coût.}
\iqfollow{Que provoque un octet corrompu~?}{Il casse toutes les valeurs suivantes, puisque chacune dépend du total courant. Les formats en colonnes écrivent une valeur complète au début de chaque bloc, si bien que le dégât reste local et que l'on peut y accéder directement.}
\end{solution}

\begin{solution}{iq:iv3:algorithms:49}
On utilise une table ordonnée par côté, du prix vers un niveau portant la quantité totale et le nombre d'ordres, ainsi qu'une table de hachage de l'identifiant vers l'ordre et un itérateur vers son niveau, comme dans le \cref{lst:iv3:al:book}. Un ajout coûte $O(\log L)$, car il peut falloir trouver ou créer le niveau. Une annulation coûte $O(1)$ en moyenne, et le meilleur prix est en $O(1)$.
\omcode{code/interviews-3/01-algorithms/cpp/book.hpp}{6}{45}{Un carnet de niveaux et un index des ordres.}\label{lst:iv3:al:book}

C'est la table des identifiants qui rend les annulations peu coûteuses. Une annulation ne désigne que l'identifiant, donc sans elle il faudrait parcourir chaque niveau. Conserver l'itérateur de la table va plus loin~: il n'y a plus de second accès par prix, puisque les itérateurs d'un \texttt{std::map} survivent aux insertions et aux suppressions d'autres clés. Effacer via un itérateur est en temps constant amorti selon la norme, et il est sûr de retirer le niveau quand son dernier ordre le quitte, de sorte qu'un niveau vide n'apparaisse jamais comme meilleur. La lecture du meilleur utilise \texttt{begin} côté vendeur et \texttt{rbegin} côté acheteur, chacun en temps constant.

Les prix et les quantités sont des entiers, car un prix en virgule flottante fait une mauvaise clé de table. Les tests exécutent 500 opérations aléatoires à la fois contre un modèle en force brute écrit en Python, y compris avec des identifiants répétés et des annulations inconnues, et aussi sous AddressSanitizer.

\iqwrong{Une file de priorité sur les prix. Le meilleur prix est rapide à lire, mais une annulation doit retirer un élément au milieu, ce qu'un tas ne sait pas faire sans suppression paresseuse. La suppression paresseuse laisse aussi des niveaux périmés à ignorer à chaque lecture du meilleur.}
\iqfollow{Que changeriez-vous pour un chemin critique (hot path)~?}{L'activité réelle se concentre près du haut du carnet. Un tableau trié à plat de niveaux, ou un tableau indexé par ticks autour du prix courant, est donc préférable à un arbre à nœuds, car il touche moins de lignes de cache. Le coût se paie à l'insertion, par la copie, ou par un recentrage lors d'un grand mouvement.}
\iqfollow{Que faut-il modifier pour l'appariement~?}{Chaque niveau doit garder ses ordres dans l'ordre d'arrivée, en général une liste doublement chaînée intrusive, de sorte qu'une exécution prenne le plus ancien et qu'une annulation retire l'élément en $O(1)$.}
\end{solution}

\begin{solution}{iq:iv3:algorithms:50}
On construit le graphe avec une arête $a \to b$ de poids $-\log r_{ab}$, puis l'on cherche un cycle négatif avec Bellman--Ford\index{Bellman--Ford}, comme dans le \cref{lst:iv3:al:fx}. Un produit de taux supérieur à 1 est une somme de $-\log$ inférieure à 0. Le coût est en $O(n \cdot E)$, soit au plus $O(n^3)$ pour un tableau complet. Examiner les triangles ne suffit pas, car une boucle rentable peut exiger davantage de sauts.
\omcode{code/interviews-3/01-algorithms/python/iv_fx_cycle.py}{4}{30}{Bellman--Ford sur les taux passés en moins logarithme.}\label{lst:iv3:al:fx}

Initialiser chaque distance à 0 revient à relier une source cachée à tous les nœuds, de sorte qu'un cycle est trouvé où qu'il se trouve. Après $n - 1$ tours, tous les plus courts chemins sans cycle négatif sont fixés, donc une modification au tour $n$ ne peut venir que d'un tel cycle. Remonter $n$ pas depuis le nœud modifié est l'étape que l'on oublie~: ce nœud peut ne faire que pendre hors du cycle, alors que $n$ pas en arrière doivent se trouver sur le cycle. De là, on avance jusqu'à revenir au point de départ.

On prend $A \to B$ à 2, $B \to C$ à 2, $C \to D$ à $\tfrac{1}{2}$, $D \to A$ à $\tfrac{3}{5}$, et les diagonales $A \to C$ à 3, $C \to A$ à $\tfrac{1}{5}$, $B \to D$ et $D \to B$ à $\tfrac{4}{5}$. Ce tableau compte quatre triangles, dont les produits valent 0.8, 0.9, 0.8 et 0.96, tandis que $A \to B \to C \to D \to A$ se multiplie en $2 \times 2 \times \tfrac{1}{2} \times \tfrac{3}{5} = 1.2$. Sur 600 tableaux aléatoires à au plus cinq devises, le résultat coïncide avec l'énumération exacte de tous les cycles en fractions, et chaque cycle renvoyé a un produit supérieur à 1.

La tolérance \texttt{eps} compte. Des taux de $1/4$, $4/9$ et $9$ se multiplient exactement à 1, mais leurs logarithmes en virgule flottante s'additionnent à environ $-4 \times 10^{-16}$, et sans tolérance cette fonction signale cette boucle comme un gain. Le prix à payer est qu'un gain inférieur à environ $10^{-9}$ par saut est ignoré. Des frais proportionnels $f$ par conversion ajoutent $-\log(1 - f)$ à chaque arête~: des frais de 10~\% annulent ce 1.2, puisque $1.2 \times 0.9^4 < 1$.

\iqwrong{Dijkstra. Sa preuve exige des arêtes de poids positifs ou nuls, si bien qu'il ne donne ici ni les plus courts chemins, ni ne détecte de cycle.}
\iqfollow{Un desk réel exécute-t-il chaque cycle trouvé~?}{Non. Les taux ne valent qu'indicativement pour une taille donnée, chaque jambe prend du temps et peut bouger, et une cotation périmée ressemble à un gain. Un cycle n'est donc qu'un candidat, à vérifier contre des prix et des tailles exécutables.}
\end{solution}

\begin{solution}{iq:iv3:algorithms:51}
On garde \texttt{tails[n]}, la plus petite valeur pouvant terminer une suite croissante de longueur $n + 1$, et pour chaque élément on retient son prédécesseur, comme dans le \cref{lst:iv3:al:lis}. Le temps est en $O(n \log n)$ et la mémoire en $O(n)$. Sur \texttt{[3, 4, 1]}, la réponse est \texttt{[3, 4]}.
\omcode{code/interviews-3/01-algorithms/python/iv_lis.py}{4}{23}{Plus petites queues, avec prédécesseurs.}\label{lst:iv3:al:lis}

Une queue plus petite n'est jamais moins bonne, car tout ce qui prolonge une suite se terminant par 5 prolonge aussi une suite de même longueur se terminant par 3. \texttt{tails} est toujours strictement croissante, donc une recherche dichotomique trouve où \texttt{x} doit se placer~: \texttt{x} prolonge la plus longue suite dont la queue lui est inférieure. Cette position donne la longueur de la meilleure suite se terminant en \texttt{x}, et \texttt{tail\_at[n - 1]} est l'élément qui le précède à cet instant. Ainsi \texttt{prev} enregistre une vraie chaîne de positions dans l'ordre d'origine, et la parcourir depuis la dernière queue donne la réponse à rebours.

Les tests comparent la longueur et la validité à une recherche sur toutes les sous-suites, sur 400 listes aléatoires.

\iqwrong{Renvoyer \texttt{tails}. Sur \texttt{[3, 4, 1]} elle vaut finalement \texttt{[1, 4]}, qui a la bonne longueur mais n'est pas une sous-suite~: le 1 vient après le 4. Remplacer une queue enregistre une meilleure option pour la suite, et non la réponse acquise jusque-là. C'est pourquoi \texttt{prev} est nécessaire.}
\iqfollow{Et pour une suite non décroissante~?}{On utilise \texttt{bisect\_right}, afin qu'une valeur égale prolonge la suite au lieu de remplacer sa queue.}
\iqfollow{Et si les boîtes $(w, h)$ ne s'emboîtent que lorsque les deux dimensions sont strictement plus grandes~?}{On trie par $w$ croissant, et à égalité de $w$ par $h$ décroissant. On lance ensuite le même balayage sur les valeurs de $h$. L'ordre décroissant en cas d'égalité empêche d'enchaîner deux boîtes de même largeur.}
\iqfollow{Pourquoi pas la programmation dynamique classique~?}{$O(n^2)$ est plus simple, et suffit jusqu'à quelques milliers d'éléments.}
\end{solution}

\begin{solution}{iq:iv3:algorithms:52}
On trie les positions selon la première liste. Les paires discordantes sont alors exactement les inversions de la seconde liste dans cet ordre, que le tri fusion compte en $O(n \log n)$, comme dans le \cref{lst:iv3:al:kendall}. Avec $P = n(n-1)/2$ paires et $D$ inversions, $\tau = (P - 2D)/P$.
\omcode{code/interviews-3/01-algorithms/python/iv_kendall.py}{4}{28}{Comptage des inversions pendant la fusion.}\label{lst:iv3:al:kendall}

La première liste étant triée, une paire est discordante exactement lorsque la seconde liste est désordonnée à ces deux positions. Pendant la fusion, lorsqu'un élément de \texttt{right} est pris avant les éléments encore présents dans \texttt{left}, il forme une inversion avec chacun d'eux d'un seul coup, et c'est ce qui rend le comptage non quadratique. Pour \texttt{[1, 2, 3, 4]} et \texttt{[1, 3, 2, 4]}, il y a 1 inversion parmi 6 paires, donc $\tau = 2/3$. Les tests concordent avec un comptage en $O(n^2)$ et avec SciPy sur 200 cas aléatoires.

En présence d'égalités, on trie par $(x, y)$, de sorte qu'une paire égale en $x$ soit ordonnée par $y$ et ne soit jamais comptée comme inversion, et qu'une paire égale en $y$ ne soit pas non plus une inversion, puisque la comparaison est $\le$. On compte ensuite $t_x$, $t_y$ et $t_{xy}$, les paires égales en $x$, en $y$ et dans les deux à la fois, et l'on utilise la forme corrigée des égalités
\[
\tau_b = \frac{C - D}{\sqrt{(P - t_x)(P - t_y)}}, \qquad C = P - t_x - t_y + t_{xy} - D.
\]
La version des tests concorde aussi avec celle par défaut de SciPy sur des données comportant des égalités.

\iqwrong{Utiliser un entier 32 bits pour les compteurs en C++. Pour $n = 100\,000$, il y a déjà 4\,999\,950\,000 paires, soit plus que $2^{31}$.}
\iqfollow{Pourquoi pas un arbre de Fenwick~?}{Cela fonctionne aussi~: on remplace les valeurs de $y$ par leur rang, puis on compte pour chacune combien des précédentes lui sont supérieures. C'est encore en $O(n \log n)$, et cela s'adapte à des paires arrivant dans l'ordre de $x$ si les valeurs de $y$ sont connues à l'avance ou tiennent dans un petit intervalle. Sinon, il faut un arbre équilibré avec des compteurs par sous-arbre.}
\end{solution}

\begin{solution}{iq:iv3:algorithms:53}
On garde jusqu'à $k - 1$ compteurs. Un symbole connu en ajoute un, un nouveau symbole prend un compteur libre, et s'il n'en reste aucun, on décrémente tous les compteurs et l'on retire ceux qui tombent à zéro, comme dans le \cref{lst:iv3:al:heavy}. C'est le résumé de Misra--Gries\index{Misra--Gries, résumé de}. Le coût est en $O(1)$ par élément, en moyenne amortie. Chaque compteur sous-estime son symbole d'au plus $n/k$. Pour $k = 2$, c'est le vote majoritaire\index{vote majoritaire}.
\omcode{code/interviews-3/01-algorithms/python/iv_heavy.py}{1}{14}{Annulation de symboles différents ensemble.}\label{lst:iv3:al:heavy}

La preuve compte les annulations. Chaque tour de décrémentation retire une unité à chacun des $k - 1$ compteurs, et jette une unité pour le symbole qui arrive. Cela fait $k$ unités retirées, toutes de symboles différents, donc il y a au plus $n / k$ tels tours. Un symbole perd au plus une unité par tour, si bien que son nombre réel est au plus son compteur plus $n / k$. Un symbole apparaissant plus de $n/k$ fois a donc encore un compteur strictement positif, et figure dans la sortie.

Le coût par élément~: une manche de décrémentation parcourt $k - 1$ compteurs, mais les manches sont au plus $n / k$, donc l'ensemble coûte $O(n)$. Les tests vérifient la borne sur 500 flux aléatoires.

Le résultat peut contenir des faux positifs. Avec $k = 2$, le flux \texttt{a, b, c} se termine avec \texttt{c} comme candidat, qui apparaît une fois sur trois. Si le flux peut être rejoué, une seconde passe compte exactement les candidats et supprime les faux.

\iqwrong{L'échantillonnage. Il donne une estimation avec une probabilité d'erreur, pas une garantie~: un symbole lourd peut être manqué. Garder un compteur par symbole est exact, mais sa mémoire croît avec le nombre de symboles, et non avec $k$.}
\iqfollow{Comment indiquez-vous l'erreur sur un compte~?}{Le compte réel se situe entre le compteur et le compteur plus $R$, le nombre de manches de décrémentation. Aucun état supplémentaire n'est nécessaire pour trouver $R$~: chaque manche retire $k$ unités, donc $R = (n - \sum \text{compteurs}) / k$, et ce nombre est souvent bien inférieur à $n / k$.}
\end{solution}

\begin{solution}{iq:iv3:algorithms:54}
Avec $k = 1$, la réponse est 5, de 2 à 7, et avec $k = 2$ elle vaut 9, en ajoutant 1 à 5. Le \Cref{lst:iv3:al:ktrades} conserve deux nombres pour chaque nombre d'allers-retours achevés, en temps $O(nk)$ et mémoire $O(k)$. Dès que $k \ge n/2$, la réponse est la somme de toutes les hausses, qui vaut ici également 9.
\omcode{code/interviews-3/01-algorithms/python/iv_k_trades.py}{1}{10}{États selon le nombre d'allers-retours achevés.}\label{lst:iv3:al:ktrades}

Après chaque jour, vous êtes à plat avec $j$ allers-retours effectués, ou vous détenez dans le $j$-ième aller-retour. Pour détenir dans le $j$-ième aller-retour, vous le déteniez déjà, ou vous achetez avec la richesse des $j - 1$ allers-retours terminés. Pour être à plat après $j$ allers-retours, vous l'étiez déjà, ou vous vendez aujourd'hui. La boucle interne n'a pas de conflit~: un achat et une vente le même jour au même prix ne rapportent rien.

Lorsque $k \ge n/2$, la limite ne peut pas être atteinte. Une liste de $n$ prix comporte au plus $n/2$ hausses distinctes, donc vous pouvez toutes les prendre, un aller-retour par suite de hausses. Ce raccourci compte~: $k$ peut valoir un milliard, et un tableau à $k$ colonnes ne tiendrait alors pas en mémoire.

Cas limites à vérifier~: $k = 0$ donne 0, et il en va de même pour une liste décroissante, car \texttt{flat} démarre à 0 et aucune transaction n'est forcée. Chaque aller-retour supplémentaire ne peut qu'aider, donc la réponse ne diminue jamais lorsque $k$ augmente.

\iqwrong{«~Prendre les $k$ plus fortes hausses.~» Sur \texttt{[1, 5, 4, 9]} avec $k = 1$, la plus forte hausse isolée entre voisins est 5, et la vraie réponse est 8, en conservant la position pendant la baisse. Il faut fusionner les hausses, et pas seulement les choisir.}
\iqfollow{Existe-t-il quelque chose de plus rapide que $O(nk)$~?}{Oui~: fusionner les suites avec un tas et une pile donne $O(n \log n)$ quel que soit $k$. Dire qu'il existe~; peu d'intervieweurs l'attendent en entier.}
\end{solution}

\begin{solution}{iq:iv3:algorithms:55}
La réponse est 8, avec les barres 4 et 5 à la hauteur 4. Le \Cref{lst:iv3:al:histogram} est une \index{pile monotone}pile monotone, en temps $O(n)$ et mémoire $O(n)$.
\enlargethispage{-3\baselineskip}%
\omcode{code/interviews-3/01-algorithms/python/iv_histogram.py}{1}{10}{Une pile de barres croissantes, avec un 0 final.}\label{lst:iv3:al:histogram}

Le plus grand rectangle a la hauteur d'une barre, la plus basse de son intervalle, et il s'étend à gauche et à droite jusqu'à ce qu'une barre plus basse l'arrête. Pour chaque barre, il faut donc trouver la barre plus basse la plus proche de chaque côté. La pile conserve des positions aux hauteurs croissantes. Lorsqu'une nouvelle barre n'est pas plus haute que le sommet, le sommet a atteint sa limite droite en $i$, et sa limite gauche est le nouveau sommet placé juste en dessous~; sa largeur totale est donc connue à cet instant. La dernière 0 vide ce qui reste.

Hauteurs égales~: le test est \texttt{>=}, donc une barre de même hauteur retire la précédente, et la largeur de cette dernière est alors calculée trop petite. Cela ne fait rien perdre~: la barre suivante de cette hauteur reste sur la pile, et sa limite gauche dépasse la barre retirée, si bien qu'elle obtient toute la largeur. Pour \texttt{[3, 3]}, la première barre calcule 3, puis la seconde calcule 6.

Chaque barre est empilée et dépilée exactement une fois, donc l'ensemble est linéaire, même si une étape peut en dépiler plusieurs.

\iqwrong{Essayer chaque paire d'extrémités et prendre la barre la plus basse. C'est correct, mais en $O(n^2)$~; un diviser pour régner sur la barre la plus basse est en $O(n \log n)$ en moyenne et en $O(n^2)$ sur une liste triée.}
\iqfollow{Et le plus grand rectangle de uns dans une matrice~?}{Traiter chaque ligne comme la base d'un histogramme, les hauteurs comptant les uns vers le haut~: $O(rc)$ au total.}
\end{solution}

\begin{solution}{iq:iv3:algorithms:56}
Après ces quatre appels, le cache contient a et c~: b n'a qu'une utilisation et a en a deux, donc b part. Le \Cref{lst:iv3:al:lfu} conserve une table de chaque fréquence vers une liste ordonnée de clés, ainsi que la plus basse fréquence utilisée, si bien que chaque opération est en $O(1)$ en moyenne.
\omcode{code/interviews-3/01-algorithms/python/iv_lfu.py}{4}{42}{Clés regroupées par nombre d'utilisations.}\label{lst:iv3:al:lfu}

Un \texttt{get} fait passer sa clé du groupe de $f$ au groupe de $f + 1$. Un \texttt{OrderedDict} retire une clé connue en $O(1)$ et conserve l'ordre d'insertion, si bien que la tête d'un groupe est l'utilisation la plus ancienne à ce compte. C'est le critère de départage que demande l'énoncé.

Pourquoi ne pas chercher le compte le plus bas~? Parce qu'il évolue d'une façon particulière. Il ne monte que d'un à la fois, lorsqu'un accès vide son groupe~: la clé est passée à $f + 1$, donc ce groupe n'est pas vide et devient le nouveau minimum. Une clé nouvelle a le compte 1, et rien ne peut être en dessous de 1, donc elle fixe \texttt{lowest} à 1 quelle que soit sa valeur précédente. Ce sont les deux seuls cas, et aucune recherche n'est jamais nécessaire, ce qui la rend constante.

L'éviction doit précéder l'insertion de la nouvelle clé, sinon la nouvelle clé pourrait s'évincer elle-même, puisqu'elle a le compte le plus bas de tous. Un groupe vidé est supprimé aussitôt. Sinon, chaque nouveau compte laisserait un groupe vide derrière lui, et une clé lue 100\,000 fois en laisserait 100\,000, dans un cache de capacité fixe.

\iqwrong{Un tas sur (compte, dernier usage). Cela coûte $O(\log n)$ par opération, ce que l'énoncé exclut. Il faut aussi une suppression paresseuse, puisque chaque accès modifie la clé de l'entrée.}
\iqfollow{Quelle est la faiblesse de la LFU~?}{Une clé très sollicitée autrefois garde un compte élevé et reste indéfiniment, tandis que les plus récentes sont évincées. Les caches réels vieillissent ou font décroître les comptes, ou utilisent une approximation par échantillonnage.}
\end{solution}

\begin{solution}{iq:iv3:algorithms:57}
Deux transactions s'affichent à 10~: 100 provenant de la première vente et 20 de la seconde. La seconde vente laisse 30 au repos, et l'achat est entièrement exécuté. Les \Cref{lst:iv3:al:matching,lst:iv3:al:matchcancel} conservent un dictionnaire qui associe à chaque prix une file de chaque côté, un tas de prix pour le meilleur, et une table qui associe chaque identifiant à son ordre. Une soumission coûte $O(t + (1 + l) \log L)$, pour $t$ transactions, $l$ niveaux atteints et $L$ niveaux au total, et une annulation coûte $O(1)$.
\omcode{code/interviews-3/01-algorithms/python/iv_matching.py}{5}{39}{Priorité prix-temps avec annulations paresseuses.}\label{lst:iv3:al:matching}
\omcode{code/interviews-3/01-algorithms/python/iv_matching.py}{41}{45}{L'annulation elle-même.}\label{lst:iv3:al:matchcancel}

La transaction se fait au prix de l'ordre au repos, 10, et non à la limite de 11 de l'achat, puisque l'ordre au repos était là en premier et proposait 10. L'achat bénéficie du meilleur prix.

Une annulation se contente de marquer l'ordre avec une quantité restante de 0 et de le retirer de la table. La file ignore ces entrées lorsqu'elles atteignent la tête. Chaque entrée est retirée une seule fois, donc cela ne coûte rien de plus en moyenne amortie, et l'on évite de chercher l'entrée au milieu d'une liste. Un niveau vidé reste dans le tas de la même façon jusqu'à atteindre le sommet, puis il sort du dictionnaire à cet instant, si bien qu'un prix n'est jamais présent deux fois dans le tas.

Les identifiants uniques de l'énoncé comptent. Un identifiant répété écraserait son entrée dans la table, et une exécution ultérieure lèverait une \texttt{KeyError}~; un moteur réel le refuse donc à l'entrée.

\iqwrong{«~Garder une seule liste triée de tous les ordres.~» Chaque insertion coûte alors $O(n)$ pour le décalage, et un carnet proche du meilleur cours est justement l'endroit où les insertions ont lieu.}
\iqfollow{Et une modification qui diminue la quantité~?}{Les bourses conservent en général la position dans la file lorsque la quantité baisse, et la perdent lorsque la quantité augmente ou que le prix change. Une telle règle relève de la plateforme elle-même, il faut donc lire sa spécification.}
\end{solution}

\begin{solution}{iq:iv3:algorithms:58}
La réponse est 12, avec $(2, 5, 7)$ et $(6, 9, 5)$. Le \Cref{lst:iv3:al:weighted} trie par heure de fin et utilise la \index{programmation dynamique}programmation dynamique avec une recherche dichotomique, en temps $O(n \log n)$ et mémoire $O(n)$.
\omcode{code/interviews-3/01-algorithms/python/iv_weighted_jobs.py}{1}{12}{Laisser la tâche, ou la prendre.}\label{lst:iv3:al:weighted}

Les tâches étant triées par heure de fin, notons \texttt{best[i]} le meilleur total obtenu avec les $i$ premières tâches. Pour la tâche $i$, il y a deux cas. La laisser de côté, et l'on garde \texttt{best[i]}. Ou la prendre, et alors seules les tâches finissant au plus tard au début de celle-ci peuvent être conservées. Celles-ci forment un préfixe de la liste triée, si bien qu'une recherche dichotomique trouve où ce préfixe s'arrête. \texttt{bisect\_right} permet précisément à une tâche de commencer au moment exact où une autre se termine.

Prendre d'abord la tâche la plus rentable choisit ici $(2, 5, 7)$ et tombe juste par chance. Elle échoue sur $(0, 10, 10)$, $(0, 4, 6)$ et $(5, 9, 6)$~: elle prend le 10, alors que les deux 6 rapportent ensemble 12.

Prendre d'abord les tâches qui finissent le plus tôt est l'algorithme glouton qui fonctionne lorsque tous les profits sont égaux, par un argument d'échange, et il est également faux ici, car sur cette liste il prend $(0, 3, 4)$ puis $(4, 7, 3)$, pour un total de 7.

\iqwrong{«~Trier par début et appeler récursivement sur la tâche suivante.~» C'est correct avec mémoïsation, mais sans mémoïsation c'est exponentiel~: chaque tâche se ramifie entre la prendre ou la laisser.}
\iqfollow{Comment retrouveriez-vous les tâches~?}{Remontez depuis la fin~: si \texttt{best[i + 1]} diffère de \texttt{best[i]}, la tâche $i$ a été prise, il faut donc sauter à son préfixe compatible.}
\end{solution}

\begin{solution}{iq:iv3:algorithms:59}
\Cref{lst:iv3:al:fenwick} est un \index{Fenwick tree}arbre de Fenwick sur les niveaux. Chaque mise à jour et chaque somme de préfixe parcourt au plus environ $\log_2 P$ entrées, et la profondeur à un prix est le total moins la somme des niveaux inférieurs. Un tableau simple donne des mises à jour en $O(1)$ et des requêtes en $O(P)$, et un tableau de totaux cumulés donne l'inverse.
\omcode{code/interviews-3/01-algorithms/python/iv_fenwick_depth.py}{1}{21}{Un arbre de Fenwick sur les niveaux de prix.}\label{lst:iv3:al:fenwick}

L'entrée $i$ de l'arbre contient la somme d'un bloc de niveaux dont la longueur est le bit à 1 de poids faible de $i$. Pour sommer un préfixe, on part de $i$ et l'on retire ce bit~: chaque étape retire un bloc et se déplace vers la gauche. Pour mettre à jour, on ajoute ce bit~: chaque étape atteint le bloc suivant qui couvre ce niveau. Dans les deux cas, $i$ a au plus $\log_2 P$ bits à traiter.

Avec 3 au prix 4, 5 au prix 7 et 2 au prix 9, \texttt{depth(7)} vaut 7, et \texttt{depth(5)} vaut également 7, car aucune taille ne se trouve en 5 ou 6, et \texttt{depth(0)} vaut 10.

Pour trouver le prix le plus élevé auquel la profondeur atteint $Q$, le côté du prix le plus haut est malaisé avec un arbre qui somme depuis le bas. On cherche donc le plus grand $p$ pour lequel \texttt{below(p)} est au plus égal au total moins $Q$, en descendant bit par bit dans l'arbre~: $O(\log P)$, sans recherche dichotomique supplémentaire, tant qu'aucune taille n'est négative.

\iqwrong{«~Utiliser un arbre de segments.~» Cela fonctionne, avec le même coût et davantage de mémoire. Il ne l'emporte que si l'on a aussi besoin de choses qu'un arbre de Fenwick ne sait pas faire, comme un minimum sur un intervalle.}
\iqfollow{Et lorsque les prix ne tiennent pas dans un petit tableau~?}{Recueillez d'abord les prix et compressez-les en rangs, si vous pouvez les voir à l'avance. Sinon, utilisez un arbre équilibré avec des sommes par sous-arbre.}
\end{solution}

\begin{solution}{iq:iv3:algorithms:60}
Il signale des occurrences aux positions 4 et 6. Le \Cref{lst:iv3:al:kmp} est l'algorithme de \index{Knuth--Morris--Pratt}Knuth--Morris--Pratt, sous sa forme en flux. Il utilise une mémoire $O(m)$ et un coût amorti $O(1)$ par élément, et il lit chaque élément une seule fois.
\omcode{code/interviews-3/01-algorithms/python/iv_kmp.py}{1}{22}{La table des échecs permet au flux de continuer.}\label{lst:iv3:al:kmp}

\texttt{k} indique combien d'éléments du motif le flux se termine maintenant. En cas de correspondance, \texttt{k} augmente. En cas de discordance, les éléments qui viennent d'être lus forment eux-mêmes un préfixe du motif, si bien qu'un préfixe plus court peut encore servir. \texttt{fail[k - 1]} est le plus long de ces préfixes, et le code reprend à partir de là.

Après une correspondance complète, \texttt{k} revient au plus long recouvrement. C'est ce qui trouve la position 6~: A, B aux positions 3 et 4 sont déjà le début de la correspondance suivante. Revenir à 0 la manque. Revenir au début est donc faux, et pas seulement lent~: sur le motif A, A, B et le flux A, A, A, B, repartir après la discordance au troisième A ne donne aucune occurrence, car à ce moment il ne contient plus qu'un A, celui du troisième élément, et il manque la vraie correspondance aux positions 2 à 4.

Pour le coût~: \texttt{k} augmente d'au plus 1 par élément, et chaque baisse le diminue d'au moins 1, sans jamais descendre sous 0, donc les baisses totalisent au plus les hausses.

\iqwrong{«~Garder les $m$ derniers éléments et comparer à chaque élément.~» C'est correct, mais en $O(m)$ par élément.}
\iqfollow{Et Rabin--Karp~?}{Garder un hachage glissant des $m$ derniers éléments~: $O(1)$ par élément, en moyenne. C'est plus simple à écrire, et il faut une comparaison en cas de coïncidence, car les hachages peuvent entrer en collision. Il gère aussi plusieurs motifs de même longueur avec un ensemble.}
\end{solution}

\begin{solution}{iq:iv3:algorithms:61}
Le plus court intervalle est A, C, B, aux positions 5 à 7. Si A est exigé deux fois, il s'allonge en A, Y, B, Z, A, C, aux positions 1 à 6. Le \Cref{lst:iv3:al:minwindow} est une \index{fenêtre glissante}fenêtre glissante, en temps $O(n + r)$ pour $r$ éléments exigés.
\omcode{code/interviews-3/01-algorithms/python/iv_min_window.py}{1}{20}{Agrandir jusqu'à la validité, rétrécir tant que la fenêtre est valide.}\label{lst:iv3:al:minwindow}

\texttt{missing[s]} est le nombre d'exemplaires de $s$ dont la fenêtre a encore besoin. Il devient négatif lorsque la fenêtre contient des exemplaires en trop, et les symboles jamais exigés partent de 0 et descendent aussitôt sous zéro. \texttt{short} compte les éléments exigés qui ne sont pas encore couverts, et la fenêtre est valide exactement lorsqu'il atteint 0. Il ne change que lorsqu'une entrée de \texttt{missing} franchit zéro, ce qui gère les répétitions~: le second A ne fait baisser \texttt{short} que si un second A a été exigé. Le listing suppose qu'au moins un symbole est exigé~: sans aucun, \texttt{short} part de 0 et le rétrécissement sort de la fin du flux.

Pour chaque extrémité droite, le code réduit la fenêtre par la gauche tant qu'elle reste valide, et retient la meilleure. Chaque position entre et sort une seule fois, donc le total est linéaire. L'extrémité droite n'avance qu'une fois la fenêtre devenue invalide, et pour cette extrémité le pointeur gauche a déjà dépassé tous les débuts qui pouvaient convenir, car la validité d'un intervalle plus long découle de celle d'un intervalle plus court.

\iqwrong{«~Suivre l'ensemble des symboles encore manquants.~» Un ensemble ne distingue pas un A de deux, alors que l'énoncé indique que les symboles exigés peuvent se répéter. Sur la seconde variante, il s'arrête à la première fenêtre d'apparence valide, qui ne contient qu'un A.}
\iqfollow{Et si chaque symbole a un prix et que vous voulez l'intervalle valide le moins cher, et non le plus court~?}{Si les prix sont positifs, la même fenêtre convient, avec un coût cumulé à la place d'une longueur, car réduire la fenêtre fait toujours baisser le coût.}
\end{solution}

\begin{solution}{iq:iv3:algorithms:62}
Pour 100 millions de clés à 1~\%, compter environ 9.6 bits par clé~: 958.5 millions de bits, soit environ 120~Mo, avec 7 fonctions de hachage. Le \Cref{lst:iv3:al:bloom} en est une version réduite. Un faux négatif est impossible, puisque l'ajout d'une clé met à 1 tous ses bits, et que les bits ne reviennent jamais à 0.
\omcode{code/interviews-3/01-algorithms/python/iv_bloom.py}{1}{26}{Un filtre de Bloom à partir d'un hachage de 128~bits.}\label{lst:iv3:al:bloom}

Un bit donné reste à 0 après $n$ insertions avec la probabilité $(1 - 1/m)^{hn} \approx e^{-hn/m}$. Un faux positif exige que les $h$ bits soient tous à 1~: $p \approx (1 - e^{-hn/m})^h$. Minimiser sur $h$ donne $h = (m/n)\ln 2$, où la moitié des bits sont à 1. En substituant, on obtient $p = 2^{-h}$, donc $m = -n \ln p / (\ln 2)^2$. Pour $p = 0.01$, cela donne $9.585\,n$, et $h = 6.64$, arrondi à 7.

Le test le mesure~: 20\,000 clés dans un tel filtre donnent un taux de faux positifs voisin de 1~\% sur 200\,000 clés nouvelles, et diviser par deux le nombre de bits le porte à environ 16~\%, ce qui correspond à la formule.

On ne peut pas supprimer, car un bit peut être partagé avec une autre clé, et le remettre à 0 créerait des faux négatifs. Un filtre à compteurs contient un petit compteur par case à la place d'un bit, pour environ quatre fois plus de mémoire.

Les deux hachages de 64 bits du listing donnent les sept cases par $a + ib$, une façon courante d'éviter de calculer sept hachages.

\iqwrong{«~Un faux positif veut dire qu'on rejette un vrai ordre.~» Seulement si on s'en sert ainsi. Il faut utiliser le filtre pour éviter le travail sur un identifiant sûrement nouveau, et confirmer un «~peut-être~» auprès de la base de référence.}
\iqfollow{Par rapport à une table de hachage~?}{Un ensemble d'identifiants de 8 octets à un facteur de charge de 0.5 occupe environ 16 octets par clé, soit 1.6~Go ici~: plus de dix fois plus, et sans aucune erreur.}
\end{solution}

\begin{solution}{iq:iv3:algorithms:63}
Le \Cref{lst:iv3:al:bestvenue} est un \index{tas indexé}tas indexé~: un tas ordinaire de plateformes, plus un tableau indiquant où se trouve chaque plateforme. \texttt{update} fait monter ou descendre son entrée en $O(\log V)$, et \texttt{best} lit la racine en $O(1)$.
\omcode{code/interviews-3/01-algorithms/python/iv_best_venue.py}{1}{29}{Un tas qui connaît la position de chaque plateforme.\protect\hfil\protect\break}\label{lst:iv3:al:bestvenue}

Un tas de bibliothèque (\texttt{heapq}, \texttt{std::priority\_queue}) n'offre que \texttt{push} et \texttt{pop}, parce qu'il ne sait pas retrouver une entrée. Modifier le prix d'une plateforme doit partir de la position de cette plateforme, et le tableau \texttt{pos} la fournit. Chaque permutation doit le mettre à jour, donc toutes les permutations passent par une seule fonction. Un cours acheteur peut bouger dans les deux sens~: une boucle remonte, la suivante descend, et au plus l'une des deux fait quelque chose.

Il existe aussi une version paresseuse~: on empile (prix, plateforme) à chaque mise à jour, et à la lecture de la racine, on dépile toute entrée dont le prix n'est plus celui actuel de la plateforme. Elle est simple, mais le tas grossit avec les mises à jour entre deux lectures, et le coût de \texttt{best} devient amorti.

Avec $V = 8$, on n'utiliserait pas de tas du tout. Un parcours linéaire de huit nombres sur une seule ligne de cache prend quelques comparaisons, sans suivre de pointeurs, et peut être vectorisé. Ou bien on garde la meilleure plateforme, et on ne refait le parcours que lorsque le cours acheteur de cette plateforme baisse. Le $\log$ n'apporte un gain qu'à partir du moment où $V$ est grand~: il faut dire ce que l'on mesurerait avant de choisir.

\iqwrong{«~Trier les plateformes à chaque mise à jour.~» Cela coûte $O(V \log V)$ à chaque tick. Déplacer un seul élément mal placé ne coûte que $O(V)$, même avec une simple insertion.}
\iqfollow{Et les égalités~?}{Il faut dire quelle règle on veut. Des cours acheteurs égaux sur deux plateformes peuvent demander une priorité selon les frais, et la comparaison porte alors sur (prix, priorité).}
\end{solution}

\begin{solution}{iq:iv3:algorithms:64}
On utilise un \index{tri par fusion externe}tri par fusion externe~: produire des séries triées qui tiennent en mémoire, puis les fusionner. Avec 30~Go de mémoire de travail, cela donne 334 séries. Une seule passe de fusion suffit, et l'ensemble du travail prend environ 22 heures de temps disque.
\enlargethispage{-3\baselineskip}%
\omcode{code/interviews-3/01-algorithms/python/iv_external_sort.py}{1}{26}{Séries, puis une fusion à k voies.}\label{lst:iv3:al:external}

\emph{Première phase.} On lit 30~Go, on trie, on écrit une série. Ainsi $10\,000 / 30$ s'arrondit à 334 séries, et la phase lit et écrit chacune 10~To.

\emph{Deuxième phase.} On fusionne avec un tampon d'entrée par série et un tampon de sortie~: un tas sur les têtes des séries. Avec, disons, 64~Mo par tampon, 334 séries occupent environ 21~Go, ce qui tient. Une seule passe suffit donc, et elle lit et écrit de nouveau 10~To. Une seconde passe ne serait nécessaire que si les séries étaient plus nombreuses que les tampons.

\emph{Temps.} Le total fait 4 passes de 10~To, à 500~Mo/s. Chacune prend $10^{13} / (5 \times 10^8) = 20\,000$ secondes, donc le total vaut $80\,000$ secondes, soit environ 22.2 heures. Le tri en mémoire prend aussi du temps, mais à cette échelle, c'est le disque qui domine.

Sur un disque à plateaux, les grands tampons comptent, car ils transforment un déplacement de tête par enregistrement en un déplacement de tête par tampon. Un tas binaire coûte au plus environ $2 \log_2 334 \approx 17$ comparaisons par enregistrement, et un arbre des perdants environ $8.4$.

\iqwrong{«~Le mapper en mémoire et appeler sort.~» Un tri par tas saute d'un endroit à l'autre du fichier, donc chaque comparaison peut provoquer un défaut de page. Même un tri rapide, dont les parcours sont séquentiels, fait environ 8 niveaux de lectures et d'écritures complètes avant qu'une partie tienne en mémoire, contre 2 pour la fusion.}
\iqfollow{Et avec beaucoup de machines~?}{Partitionner par intervalles de clés, choisis à partir d'un échantillon pour que les parts soient égales, envoyer chaque enregistrement à sa machine et l'y trier. Le résultat est trié d'abord par machine, puis à l'intérieur de chacune.}
\end{solution}

\begin{solution}{iq:iv3:algorithms:65}
Le \Cref{lst:iv3:al:alias} est la \index{méthode d'alias}méthode d'alias. Elle répartit la probabilité en $n$ cases de même poids total, chacune contenant au plus deux éléments. Un tirage choisit une case uniformément, puis lance une pièce entre les deux éléments. Pour les poids donnés, la case 0 ne contient que l'élément 0, la case 1 seulement l'élément 1, et la case 2 contient l'élément 2 avec probabilité $\frac{1}{2}$, et l'élément 0 sinon.
\omcode{code/interviews-3/01-algorithms/python/iv_alias.py}{1}{20}{Construction de la table et tirage.}\label{lst:iv3:al:alias}

On met les poids à l'échelle en les multipliant par $n$, de sorte que la moyenne vaille 1~: ici $\frac{3}{2}$, 1 et $\frac{1}{2}$. Tout élément inférieur à 1 remplit en partie sa propre case et prend le reste à un élément supérieur à 1, qui perd alors cette quantité. Le procédé se termine toujours~: chaque étape remplit une case, et si un élément est inférieur à 1, un autre doit être supérieur à 1, puisque la moyenne vaut exactement 1.

Ainsi l'élément 2 (mis à l'échelle à $\frac{1}{2}$) remplit la moitié de la case 2, et l'élément 0 remplit le reste, ce qui le ramène à 1 et complète sa propre case.

Vérifions le résultat. L'élément 0 est tiré avec la probabilité $\frac{1}{3} + \frac{1}{3} \times \frac{1}{2} = \frac{1}{2}$. L'élément 1 reçoit $\frac{1}{3}$. L'élément 2 reçoit $\frac{1}{3} \times \frac{1}{2} = \frac{1}{6}$. Tout est juste, et dans le test, une somme exacte et une simulation concordent. Cette table provient de fractions exactes, comme dans le test. Avec des flottants, $\frac{1}{3} \times 3$ divisé par le total peut tomber juste au-dessus de 1, et le listing construit alors une autre table valide, aux mêmes probabilités.

\iqwrong{«~Construire un tableau d'un million d'entrées proportionnel aux poids.~» Le tirage coûte alors $O(1)$, mais l'arrondi déforme les petits poids, et la mémoire dépend de la précision souhaitée.}
\iqfollow{Et si les poids changent~?}{La mise en place coûte à nouveau $O(n)$. Si les poids changent souvent, il faut un arbre de sommes partielles, en $O(\log n)$ pour les deux opérations.}
\end{solution}

\begin{solution}{iq:iv3:algorithms:66}
Il n'y a aucune action devant le troisième ordre, qui en a encore 250~: la transaction prend les 100 du premier ordre puis, le deuxième ayant disparu, 50 du troisième lui-même. Le \Cref{lst:iv3:al:queuepos} conserve un \index{arbre de Fenwick}arbre de Fenwick sur les numéros d'arrivée, qui contient les actions au repos.
\omcode{code/interviews-3/01-algorithms/python/iv_queue_position.py}{1}{39}{Un arbre de Fenwick sur l'ordre d'arrivée.}\label{lst:iv3:al:queuepos}

On attribue à chaque ordre le numéro d'arrivée suivant. L'arbre enregistre les actions encore au repos à chaque numéro. Une annulation les retranche, et le nombre d'actions devant un ordre est une somme préfixe sur les numéros antérieurs. Chacune coûte $O(\log n)$.

Une transaction parcourt la file depuis l'avant, en réduisant chaque ordre à son tour. Elle peut toucher beaucoup d'ordres, mais chaque ordre n'est vidé qu'une fois au plus, donc le coût amorti est de $O(\log n)$ par ordre consommé, plus un pour l'exécution partielle.

Un total cumulé attaché à chaque ordre ne tient pas face aux annulations. Si un ordre placé devant lui quitte la file, chaque ordre situé derrière doit être mis à jour, ce qui coûte $O(n)$. L'arbre reporte cette somme jusqu'à ce que quelqu'un la demande.

Les numéros d'arrivée croissent sans fin, donc une longue journée demande une borne ou une reconstruction périodique qui compacte les numéros. La borne du tableau est ici \texttt{max\_orders}, et le test la vérifie contre une simple liste sur des événements aléatoires.

\iqwrong{«~Le nombre d'actions devant un ordre est la somme des ordres qui le précèdent, donc il suffit de parcourir la liste.~» C'est juste, à $O(n)$ par question, et elle revient en général à chaque mise à jour.}
\iqfollow{Pourquoi voudriez-vous ce nombre~?}{L'estimation que fait un teneur de marché de sa chance d'être exécuté avant que le prix ne bouge en dépend. Elle est calculée à presque chaque mise à jour, pour chaque ordre au repos.}
\end{solution}

\begin{solution}{iq:iv3:algorithms:67}
Le \Cref{lst:iv3:al:snapshot} comporte deux états~: tant qu'il n'est pas synchronisé, il met les mises à jour en tampon. Quand un instantané arrive, il supprime les mises à jour mises en tampon dont le numéro de séquence est inférieur ou égal à celui de l'instantané, applique le reste dans l'ordre des numéros de séquence, puis passe à l'état synchronisé. Une fois synchronisé, il applique le numéro suivant et traite tout autre numéro comme un doublon ou un trou, et il ignore les instantanés, qui ne pourraient que faire revenir le carnet en arrière.
\omcode{code/interviews-3/01-algorithms/python/iv_snapshot.py}{1}{39}{Mise en tampon, application d'un instantané, rejeu.}\label{lst:iv3:al:snapshot}

La mise à jour qui arrive pendant l'attente est conservée, car l'instantané peut être plus ancien qu'elle. La rejeter et repartir de l'instantané perdrait chaque changement survenu entre les deux. Une mise à jour de numéro inférieur ou égal à celui de l'instantané est déjà incluse dans celui-ci~; l'appliquer rejouerait donc un état ancien~: comme chaque mise à jour fixe une taille, cela pourrait rétablir une taille périmée.

L'instantané doit aussi se raccorder~: les mises à jour mises en tampon au-dessus de lui doivent être consécutives à partir de son numéro plus un. S'il est plus ancien que la première mise à jour mise en tampon, il y a un trou, et \texttt{on\_snapshot} renvoie faux et continue d'attendre le suivant. C'est pourquoi le tampon est conservé entre deux tentatives échouées, et il lui faut une borne dans le code réel.

En cas de trou alors que le carnet est synchronisé, celui-ci est jeté aussitôt, pour que personne ne négocie dessus, et la mise à jour qui arrive ouvre le tampon. Un doublon issu de canaux redondants, qui arrivent souvent par paires de flux, est ignoré.

Un contrôle aux limites~: sans perte, le carnet ne quitte jamais l'état synchronisé et chaque mise à jour s'applique à son tour. Dans le test, les pertes et les réordonnancements aléatoires, comparés à un modèle, aboutissent toujours au carnet du modèle, ou à un état déclaré non synchronisé, et jamais à un carnet faux.

\iqwrong{«~Appliquer l'instantané, puis appliquer tout le tampon.~» Les mises à jour périmées écrasent un état plus récent, et le carnet paraît plausible, sans erreur, jusqu'à ce qu'une transaction s'y appuie.}
\iqfollow{Et si la taille d'un niveau de prix est transmise comme une variation et non comme une nouvelle valeur~?}{La même logique s'applique, et une mise à jour appliquée deux fois donne alors une taille fausse~: un doublon est donc plus dangereux.}
\end{solution}

\begin{solution}{iq:iv3:algorithms:68}
Le coût minimal est 33, en fusionnant 2 et 3 pour 5, puis 4 et 5 pour 9, et enfin 9 et 10 pour 19. En partant du plus grand, le coût est 50~: 14, puis 17, puis 19. Le \Cref{lst:iv3:al:mergeruns} fusionne toujours les deux plus petites séries, ce qui est la \index{codage de Huffman}règle de Huffman, en temps $O(n \log n)$.
\omcode{code/interviews-3/01-algorithms/python/iv_merge_runs.py}{1}{25}{Toujours fusionner les deux plus petites séries.}\label{lst:iv3:al:mergeruns}

Les fusions forment un arbre binaire dont les feuilles sont les séries, et la taille de chaque série est comptée une fois pour chaque fusion située au-dessus d'elle. Le coût total est donc la somme des tailles multipliées par les profondeurs, et l'objectif est de placer les grandes séries près de la racine.

La démonstration repose sur un argument d'échange. Dans un arbre optimal, deux feuilles du niveau le plus profond sont des feuilles sœurs. Si les deux plus petites séries n'étaient pas au niveau le plus profond, les échanger avec ce qui s'y trouve ne pourrait pas augmenter le coût, puisque la taille la plus grande remonte. Il existe donc un arbre optimal dont les deux plus petites séries sont les feuilles sœurs les plus profondes. Les fusionner donne un problème plus petit d'une série de moins, dont le coût ne diffère que d'une quantité fixe, et l'on conclut par récurrence sur le nombre de séries.

Pour une fusion de trois à la fois, chaque fusion retire deux séries. Il ne reste donc une seule série que si $n - 1$ est pair. Sinon, ajouter d'abord des séries de taille 0. Pour les mêmes séries, $n - 1 = 3$ est impair~: on ajoute donc un 0 et on prend à chaque fois les trois plus petites~: $0 + 2 + 3 = 5$, puis $4 + 5 + 10 = 19$, soit un total de 24. Sans le 0, la première fusion devrait porter sur des séries plus grandes.

\iqwrong{«~Fusionner dans n'importe quel ordre coûte le même prix, puisque de toute façon toutes les données sont fusionnées.~» Cela vaut pour les octets de la série finale, mais pas pour le travail~: 33 contre 50.}
\end{solution}

\begin{solution}{iq:iv3:algorithms:69}
Utiliser deux tas avec suppression paresseuse, en marquant chaque entrée de sa position, comme dans le \cref{lst:iv3:al:winmed}. Cela prend un temps $O(n \log n)$ au pire. La mémoire est $O(n)$ au pire, et non $O(k)$, car les entrées qui ont quitté la fenêtre restent dans le tas jusqu'à atteindre sa racine.
\omcode{code/interviews-3/01-algorithms/python/iv_window_median.py}{4}{33}{Deux tas avec les positions dans leurs entrées.\protect\hfil\protect\break}\label{lst:iv3:al:winmed}

La structure est celle de la médiane glissante~: \texttt{low} contient la moitié inférieure, et la médiane se lit aux racines. Un tas ne permet pas de supprimer au milieu, donc une ancienne valeur reste en place et n'est écartée que lorsqu'elle atteint une racine. Marquer les entrées par leurs positions rend ce procédé sûr~: le test de péremption est \texttt{pos <= i - k}, de sorte que les doublons ne demandent aucun traitement particulier, alors qu'un compteur commun de valeurs à supprimer devrait décider de la copie du tas qu'il visait. Le compteur \texttt{live} indique combien d'entrées de \texttt{low} sont réellement dans la fenêtre, car \texttt{len(low)} ne le dit plus, et \texttt{in\_low} indique de quel côté se trouvait la valeur qui sort. Après chaque pas, au plus deux entrées changent de côté, donc le nombre de déplacements par pas est borné.

Les tas peuvent contenir presque tout ce qui a été vu. Sur une série croissante, les anciennes valeurs sont petites~; elles descendent au fond du tas-max et n'en remontent jamais~: l'exécution du test avec $n = 10\,000$ et $k = 5$ a gardé presque les $10\,000$ entrées à son maximum. Chaque opération sur un tas coûte donc $O(\log n)$, et non $O(\log k)$, et sur un flux sans fin la mémoire croît sans limite. Pour y remédier, reconstruire les tas à partir de la fenêtre dès que les entrées périmées sont plus nombreuses que les entrées vivantes, ce qui garde la mémoire à $O(k)$ et le temps par pas à $O(\log k)$ amorti.

\iqwrong{Annoncer $O(n \log k)$ sans reconstruction. La borne séduisante suppose que les tas ne contiennent que la fenêtre.}
\iqfollow{Qu'est-ce qui donnerait $O(\log k)$ directement~?}{Un arbre équilibré avec suppression réelle, ou deux ensembles triés. Ou un arbre de Fenwick sur des valeurs recodées par leur rang, avec une recherche du rang médian~: cela coûte $O(n \log n)$ sans entrées périmées, mais hors ligne.}
\iqfollow{Quelle est la solution la plus simple en pratique~?}{Un tableau trié avec recherche dichotomique et déplacement mémoire, qui coûte $O(k)$ par pas, avec une petite constante. Mesurer avant de l'écarter pour les petits $k$.}
\end{solution}

\begin{solution}{iq:iv3:algorithms:70}
Calculer les sommes préfixes $S_0, \dots, S_n$, et parcourir les extrémités $j$ avec une file double des débuts candidats dont les valeurs de $S$ croissent strictement, comme dans le \cref{lst:iv3:al:shortest}. Chaque indice entre et sort au plus une fois, donc le temps est $O(n)$. L'exemple donne 3~: le segment \texttt{32, 40, 95} a pour somme 167.
\omcode{code/interviews-3/01-algorithms/python/iv_shortest_sum.py}{4}{19}{Une file double monotone sur les sommes préfixes.}\label{lst:iv3:al:shortest}

On veut $j - i$ aussi petit que possible avec $S_j - S_i \ge T$. Deux règles élaguent les débuts, et chacune est sûre pour sa propre raison.

La première règle retire par l'arrière. Si $i < i'$ et $S_i \ge S_{i'}$, alors $i$ ne peut plus jamais être le meilleur début~: pour tout $j$ ultérieur, $i'$ donne une somme au moins aussi grande et un segment plus court. Ainsi, lorsque $j$ arrive, il retire tous les candidats dont le préfixe est au moins aussi élevé.

La seconde règle retire par l'avant. Si $S_j - S_i \ge T$, alors $i$ a trouvé sa meilleure fin, puisque toute fin ultérieure donne un segment plus long à partir du même début. On enregistre donc $j - i$ et on retire $i$. Comme $T$ est positif, un début n'est jamais associé qu'à une fin postérieure.

Ensemble, ces règles maintiennent la file croissante, de sorte que si l'avant ne convient pas, rien derrière lui ne convient, et la boucle \texttt{while} peut s'arrêter au premier échec. Les tests comparent avec une recherche complète sur 800 listes aléatoires contenant des valeurs négatives.

\iqwrong{La fenêtre glissante qui se réduit par la gauche tant que la somme vaut au moins $T$. Elle n'est juste que si aucune valeur n'est négative. Sur cet exemple, elle n'atteint jamais la réponse 3, car une valeur négative rompt la règle selon laquelle réduire la fenêtre fait baisser la somme.}
\iqfollow{Et si la seule chose dont vous vous souvenez est un tas de sommes préfixes~?}{Dépilez chaque début qui vérifie l'inégalité à chaque $j$~: cela coûte $O(n \log n)$, et c'est correct. Dites-le, puis précisez que la file double améliore cette solution.}
\end{solution}

\begin{solution}{iq:iv3:algorithms:71}
Appliquer une recherche dichotomique sur la plus courte liste pour savoir combien de ses valeurs se trouvent à gauche de la coupure médiane, comme dans le \cref{lst:iv3:al:twomed}. Cela prend un temps $O(\log \min(m, n))$ et une mémoire $O(1)$. Avec un billion d'entiers répartis sur de nombreuses machines, effectuer plutôt la dichotomie sur la valeur, en demandant à chaque machine combien de ses valeurs sont au plus égales au candidat.
\omcode{code/interviews-3/01-algorithms/python/iv_two_medians.py}{5}{25}{Recherche dichotomique sur la coupure.}\label{lst:iv3:al:twomed}

Répartir toutes les valeurs entre une partie gauche de taille $\lceil (m + n)/2 \rceil$ et une partie droite. Si $i$ viennent de $a$, alors $j = \text{half} - i$ viennent de $b$. Chaque liste est triée, donc la coupure est juste exactement lorsque $a_{i-1} \le b_j$ et $b_{j-1} \le a_i$. Si $a_{i-1} > b_j$, on a pris trop d'éléments dans $a$, et l'autre échec signifie qu'on en a pris trop peu~; la condition indique donc dans quel sens se déplacer, et la médiane se trouve à la frontière. Les infinis gèrent une coupure à l'extrême bout d'une liste, et la recherche sur la plus courte liste garantit que $j$ reste dans les bornes. Le test exécute 1000 paires aléatoires, avec listes vides et doublons, contre la réponse triée.

Pour de nombreuses machines, l'astuce des deux listes ne se généralise pas proprement. Mais une médiane admet une caractérisation plus simple~: c'est la plus petite valeur $v$ telle qu'au moins la moitié de tous les éléments soient au plus égaux à $v$. On effectue donc une dichotomie sur $v$ entre 0 et $10^9$. Chaque essai demande à chaque machine un décompte, par recherche dichotomique sur ses données triées. Cela représente environ 30 tours de messages, un par bit de l'intervalle, et chaque tour ne déplace qu'un seul nombre par machine, et non les données.

\iqwrong{Faire la moyenne des médianes des deux listes. Pour \texttt{[1, 2, 3, 4, 5, 6, 7]} et \texttt{[100]}, cela donne 52, alors que la vraie médiane vaut 4.5. La version classique, qui compare les médianes et écarte des moitiés, ne fonctionne qu'avec des longueurs égales et des bords soigneusement traités.}
\iqfollow{Le $k$-ième plus petit élément des deux listes~?}{Comparer les éléments vers la position $k/2$ dans chaque liste, et écarter le préfixe qui se termine par le plus petit, puisque tout ce préfixe est situé sous le $k$-ième. Cela demande environ $\log k$ étapes.}
\end{solution}

\begin{solution}{iq:iv3:algorithms:72}
Il renvoie~: à l'instant 0 le prix 10, à 1 le prix 12, à 3 le prix 11, à 4 le prix 10, à 5 rien, à 6 le prix 9, et à 7 rien. Le \Cref{lst:iv3:al:bidhistory} balaie les instants d'événements avec un tas des cours acheteurs actifs, en ne retirant les cours périmés que lorsqu'ils atteignent le sommet, en temps $O(n \log n)$.
\enlargethispage{-3\baselineskip}%
\omcode{code/interviews-3/01-algorithms/python/iv_best_bid_history.py}{1}{18}{Un balayage avec un tas paresseux.}\label{lst:iv3:al:bidhistory}

Le meilleur cours acheteur ne peut changer qu'à une ouverture ou à une clôture, de sorte que seuls les instants d'événements doivent être visités. À chaque instant, ajouter d'abord les cours acheteurs qui s'ouvrent alors, puis dépiler tant que le sommet est clos. Un cours acheteur périmé plus bas peut rester en place, car il ne peut affecter la réponse qu'en atteignant le sommet. Chaque cours acheteur est donc empilé une fois et dépilé au plus une fois, et le total est $O(n \log n)$.

L'ordre à l'intérieur d'un même instant compte. Sur un intervalle semi-ouvert, un cours acheteur qui se clôt à $t$ n'est plus là à $t$, et un qui s'ouvre à $t$ y est~; le code obtient ce résultat en testant \texttt{close <= t}. Si un cours acheteur se clôt à 5 et un autre s'ouvre à 5 au même prix, aucun changement n'est affiché, car la comparaison avec la dernière valeur le supprime. C'est pourquoi on compare avant d'ajouter~: l'historique contiendrait sinon un faux changement.

C'est le problème de la silhouette, sous les habits d'un livre, et la même méthode le résout. Un test contre la définition directe, en vérifiant chaque instant entier, s'exécute sur des cours acheteurs aléatoires.

\iqwrong{«~Triez par prix, et peignez chaque cours acheteur sur la durée.~» Peindre coûte $O(n \cdot \text{span})$, et ce n'est linéaire en l'étendue que si les instants sont de petits entiers.}
\iqfollow{Les cours acheteurs peuvent-ils être annulés sur un flux en direct, sans heure de clôture connue~?}{Le même tas paresseux convient, avec un ensemble d'identifiants annulés, et chaque annulation marque son entrée. À tout moment, le meilleur cours actif est le premier sommet qui n'est pas annulé.}
\end{solution}

\begin{solution}{iq:iv3:algorithms:73}
Le décompte est 5. Le \Cref{lst:iv3:al:rangesums} applique un tri fusion aux sommes préfixes, en temps $O(n \log n)$ et mémoire $O(n)$~: un segment allant de $i$ à $j$ a une somme dans l'intervalle exactement lorsque $\text{prefix}_j - \text{prefix}_i$ l'est, pour $i < j$.
\omcode{code/interviews-3/01-algorithms/python/iv_range_sums.py}{1}{21}{Décompte à travers les deux moitiés d'un tri fusion.\protect\hfil\protect\break}\label{lst:iv3:al:rangesums}

Les préfixes valent 0, 3, 1, 2, 6. Les segments dont les sommes vont de 1 à 3 sont \texttt{[3]}, \texttt{[3, -2]}, \texttt{[3, -2, 1]}, \texttt{[-2, 1, 4]} et \texttt{[1]}, et les cinq autres sortent de l'intervalle.

Une table de hachage répond à «~combien de préfixes antérieurs sont égaux à cette valeur~?~», ce qui est la question de la somme exacte. Ici, il faut savoir combien tombent dans un intervalle, et une table de hachage n'a pas d'ordre~: chaque requête devrait donc parcourir l'intervalle des valeurs possibles.

Séparer les préfixes en une moitié gauche et une moitié droite. Les paires situées à l'intérieur d'une moitié sont comptées par la récursion. Pour une paire à cheval sur les deux, les deux moitiés sont triées. Lorsque \texttt{p} parcourt la moitié gauche, la fenêtre des valeurs de droite telles que $\text{right} - p \in [L, U]$ remonte, de sorte que les deux pointeurs n'avancent que vers l'avant~: un total de $O(n)$ par niveau. Le \texttt{sorted} au retour est une fusion déguisée~: le tri de Python repère les deux suites croissantes et les fusionne, il effectue donc $O(n)$ comparaisons par niveau. Dans un autre langage, écrire la fusion.

\iqwrong{Une fenêtre glissante. Elle repose sur la croissance des sommes avec la fenêtre, et les valeurs négatives brisent cette propriété.}
\iqfollow{Une autre voie~?}{Recoder les valeurs préfixes par leur rang, et balayer avec un arbre de Fenwick des uns vus, en comptant ceux qui sont dans $[\text{prefix} - U, \text{prefix} - L]$. Le coût est le même, mais il faut connaître tous les préfixes à l'avance. Sur un vrai flux, utiliser un arbre équilibré avec des compteurs de sous-arbres.}
\end{solution}

\begin{solution}{iq:iv3:algorithms:74}
Le répartir en 3 et 2 parts, pour un coût de 22~: la première période coûte $1 \times 3 + 3^2 = 12$ et la seconde $3 \times 2 + 2^2 = 10$. Le \Cref{lst:iv3:al:slicing} achète une part à la fois, toujours dans la période où la part suivante est la moins chère.
\omcode{code/interviews-3/01-algorithms/python/iv_slicing.py}{1}{18}{La part suivante la moins chère, grâce à un tas.}\label{lst:iv3:al:slicing}

Le coût de la $(q + 1)$-ième part dans la période $t$ est le coût marginal $a_t + b_t(2q + 1)$. Il ne baisse jamais lorsque $q$ croît, puisque $b_t \ge 0$. Ainsi, dans chaque période, les parts les moins chères sont les premières, et un ensemble de $X$ parts a un coût qui est simplement la somme de ses coûts marginaux. C'est exactement le problème de choisir les $X$ entrées les moins chères parmi $T$ listes croissantes, que la fusion par tas résout en $O(X \log T)$.

Toute répartition correspond à un préfixe de chaque liste, et les $X$ plus petites entrées globales forment toujours des préfixes, car chaque liste croît. Aucune répartition ne peut donc les battre.

Vérifications aux limites~: avec $b = 0$, chaque coût marginal est constant, et toutes les parts vont à la période la moins chère. Lorsque $b_t$ devient grand dans une période, celle-ci n'en reçoit presque rien. Pour un grand $X$, les parts tendent à égaliser les coûts marginaux, ce qui est la condition du premier ordre continue.

Un coût fixe brise la méthode. Le coût n'est plus convexe en $q$~: la première part porte le coût fixe, et les suivantes non, de sorte que les coûts marginaux baissent et que l'argument du préfixe échoue. Avec un coût fixe de 100 dans chaque période, la méthode ci-dessus répartit les parts et coûte 200, alors qu'acheter tout dans une seule période ne coûte 100 qu'une fois. Utiliser la programmation dynamique sur (période, parts restantes), en $O(TX^2)$, ou essayer chaque ensemble de périodes actives.

\iqwrong{«~Répartir également.~» Cela ignore les différences entre les périodes, et ce n'est juste que lorsque celles-ci se compensent par hasard.}
\end{solution}

\begin{solution}{iq:iv3:algorithms:75}
Associer d'abord l'acheteur 1 à A laisse l'acheteur 2 sans partenaire, pour une seule paire, alors que deux sont possibles~: l'acheteur 1 avec B et l'acheteur 2 avec A. Le \Cref{lst:iv3:al:pairing} est la méthode des chemins augmentants de \index{chemin augmentant}Kuhn~: lorsque l'acheteur 2 constate que A est pris, il demande si le partenaire de A, l'acheteur 1, peut passer à quelqu'un d'autre. Elle coûte $O(VE)$.
\omcode{code/interviews-3/01-algorithms/python/iv_pairing.py}{1}{14}{Déplacer les partenaires le long d'un chemin augmentant.\protect\hfil\protect\break}\label{lst:iv3:al:pairing}

Un chemin augmentant part d'un acheteur non apparié vers un vendeur non apparié, en alternant entre des arêtes hors de l'appariement et des arêtes qui y figurent. Retourner toutes ses arêtes ajoute une paire et n'en retire aucune. Le théorème de Berge affirme qu'un appariement est maximal exactement lorsqu'aucun chemin de ce type n'existe, donc essayer chaque acheteur à tour de rôle suffit. L'ensemble \texttt{seen} maintient chaque recherche en $O(E)$, et il y a au plus $V$ recherches.

Un acheteur qui échoue une fois n'a pas besoin d'être réessayé, car les augmentations ultérieures n'ouvrent jamais de chemin qui n'existait pas. C'est pourquoi un seul passage sur les acheteurs suffit.

La recherche s'appelle récursivement une fois par arête du chemin, donc en Python un chemin de plus d'environ 1000 arêtes lève \texttt{RecursionError}. Sur de telles données, l'écrire avec une pile explicite, ou augmenter la limite.

\iqwrong{«~Glouton~: associer chaque acheteur à son premier partenaire libre.~» L'exemple montre qu'il peut s'arrêter à la moitié de ce qui est possible. Il atteint toujours au moins la moitié, car chaque paire qu'il forme peut bloquer au plus deux paires de l'appariement optimal.}
\iqfollow{Peut-on faire mieux que $O(VE)$~?}{Hopcroft--Karp trouve de nombreux chemins augmentants les plus courts à chaque phase, pour $O(E\sqrt{V})$. C'est aussi un problème de flot maximal, à capacités unitaires~: ajouter une source aux acheteurs et un puits aux vendeurs.}
\iqfollow{Et si chaque paire a aussi une valeur~?}{C'est le problème d'affectation, que la méthode hongroise résout en $O(V^3)$.}
\iqfollow{Le plus grand nombre de transactions n'est pas le plus grand volume, n'est-ce pas~?}{Exact~: avec des tailles, cela devient un problème de flot avec capacités, et l'argument des capacités unitaires cède la place à un algorithme de flot maximal.}
\end{solution}

\begin{solution}{iq:iv3:algorithms:76}
\Cref{lst:iv3:al:regex} remplit une table indexée par les positions dans le texte et dans le motif, en temps et mémoire $O(mn)$. \texttt{a*b*} ne reconnaît pas \texttt{ba}~: le \texttt{b} vient avant le \texttt{a}, et le motif n'autorise que l'autre ordre. Un retour arrière naïf sur \texttt{a*a*a*a*b} essaie toutes les façons de répartir une suite de $n$ lettres entre les quatre étoiles, soit environ $\binom{n + 4}{4}$ d'entre elles, et $s$ étoiles donnent une puissance de $n$ voisine de $s$~: c'est exponentiel en la longueur du motif.
\omcode{code/interviews-3/01-algorithms/python/iv_regex.py}{1}{13}{Une table indexée par les positions dans le texte et le motif.}\label{lst:iv3:al:regex}

Soit \texttt{ok[i][j]} l'indication que le texte à partir de $i$ reconnaît le motif à partir de $j$. Si l'élément suivant du motif porte une étoile, il y a deux choix~: l'utiliser zéro fois, en avançant dans le motif~; ou, si le caractère courant correspond, l'utiliser une fois de plus et rester dessus. Sinon, les caractères doivent correspondre, et les deux avancent. Le seul cas de base est que le texte vide reconnaît le motif vide~; un motif comme \texttt{a*} qui reconnaît le texte vide est trouvé par la règle des zéro fois.

Il n'y a que $(n + 1)(m + 1)$ états, et chacun est calculé une seule fois à partir de ses voisins, ce qui fait tout le gain par rapport au retour arrière~: les nombreuses façons de répartir la suite de \texttt{a} passent toutes par les mêmes états.

Dans le test, la table est comparée à \texttt{re.fullmatch} de Python lui-même, sur des motifs et des textes aléatoires, et à une version comptant les appels du retour arrière naïf. Sur \texttt{a*a*a*a*b} et 20 lettres \texttt{a}, le retour arrière effectue plus de 10\,626 appels avant d'échouer, alors que la table ne compte que $21 \times 10 = 210$ cellules.

\iqwrong{«~Glouton~: que chaque étoile prenne le plus possible.~» Sur \texttt{a*a} face au texte \texttt{a}, l'étoile prend le \texttt{a}, il ne reste plus rien pour le second \texttt{a}, et la réponse est donc à tort négative.}
\iqfollow{Pourquoi cela compte-t-il en dehors d'un entretien~?}{Les moteurs d'expressions régulières à retour arrière mettent un temps exponentiel sur de tels motifs. Des motifs appliqués à un texte non fiable peuvent bloquer un service, ce qui explique l'existence de moteurs en temps linéaire.}
\end{solution}

\begin{solution}{iq:iv3:algorithms:77}
La réponse est 3~: un virement règle la paire $5, -5$, et deux virements règlent le groupe $3, -1, -2$. Le \Cref{lst:iv3:al:settle} est une \index{programmation dynamique sur masques binaires}programmation dynamique sur masques binaires, en temps $O(2^n n)$ et en mémoire $O(2^n)$, sur les $n$ soldes non nuls.
\omcode{code/interviews-3/01-algorithms/python/iv_settle.py}{1}{13}{Le plus grand nombre de groupes à somme nulle, sur tous les sous-ensembles.}\label{lst:iv3:al:settle}

Tout ensemble de virements relie les contreparties en groupes connexes, et les soldes de chaque groupe s'annulent. Un groupe de taille $g$ demande au moins $g - 1$ virements, car il doit être connexe, et $g - 1$ suffisent~: on les enchaîne en reportant le solde courant. Le total est donc $n$ moins le nombre de groupes, et le but est de répartir les soldes non nuls en autant de groupes à somme nulle que possible.

Pour chaque sous-ensemble, \texttt{groups[mask]} est le plus grand nombre de groupes à somme nulle qu'il contient. On retire un élément quelconque et on ajoute 1 si l'ensemble entier a une somme nulle. L'ordre à l'intérieur du sous-ensemble est libre, si bien que cette programmation dynamique est la méthode classique pour une partition ordonnée~: en parcourant un bon ordre, la somme courante revient à 0 une fois par groupe, et chaque retour compte un groupe.

Avec 20 contreparties, cela fait environ 20 millions d'étapes, ce qui est faisable. Au-delà, le problème contient le problème de la somme d'un sous-ensemble, puisque demander si deux groupes existent revient exactement à cela~; il est donc NP-difficile, et il faut recourir à des heuristiques.

\iqwrong{«~Apparier le plus gros créancier avec le plus gros débiteur.~» Sur \texttt{[8, 7, 2, -8, -9]}, on apparie 8 avec $-9$, il reste $-1$, puis 7 avec $-8$, il reste $-1$, et 2 doit alors payer les deux~: quatre virements. Régler 8 avec $-8$ et laisser 7 et 2 payer le $-9$ en prend trois. La règle ne cherche jamais les groupes dont la somme est nulle, et sur une liste sans un tel groupe, comme \texttt{[6, 5, -4, -4, -3]}, elle trouve tout de même la meilleure réponse, 4.}
\end{solution}

\begin{solution}{iq:iv3:algorithms:78}
La réponse est 24, obtenue à partir des valeurs 7, $-2$, 4 et 15. Le \Cref{lst:iv3:al:gapped} est une programmation dynamique avec une \index{file monotone}file monotone, en temps $O(n)$ et en mémoire $O(n)$.
\omcode{code/interviews-3/01-algorithms/python/iv_gapped.py}{1}{15}{Le maximum sur la fenêtre, à l'intérieur d'une programmation dynamique.}\label{lst:iv3:al:gapped}

Soit $d_i$ la meilleure somme d'un choix se terminant en $i$. Alors $d_i = x_i + \max(0, d_{i-1}, \dots, d_{i-k})$~: soit le choix commence à cette position, soit il suit l'une des $k$ dernières positions. Le $\max$ sur une fenêtre glissante coûterait $O(k)$ à chaque pas, soit $O(nk)$ au total.

Ce maximum sur la fenêtre est exactement ce qu'une file monotone donne en $O(1)$ amorti~: des positions dont les $d$ décroissent, la tête étant la plus grande. Le facteur $k$ disparaît donc.

Dans l'exemple, $d_0 = 7$, $d_1 = 5$, $d_2 = -2$, car $-9$ suit 7 à l'écart de 2. Puis $d_3 = 4 + \max(0, 5, -2) = 9$, et $d_4 = 15 + \max(0, -2, 9) = 24$. La réponse est le plus grand $d$. Notons qu'il doit être non vide~: sur une liste entièrement négative, la réponse est donc la plus grande valeur isolée. Le zéro dans la récurrence permet à un choix de commencer n'importe où, et pourtant chaque $d_i$ inclut $x_i$.

Cas limites~: avec $k = 1$, le choix doit être contigu, on obtient donc la plus grande somme d'un segment. Avec $k \ge n$, les écarts ne comptent plus, c'est donc la somme des valeurs positives, ou la plus grande valeur isolée si aucune n'est positive.

\iqwrong{«~Prendre toutes les valeurs positives.~» Cela ignore la règle d'écart, et donne ici 26 à partir de 7, 4 et 15, alors que l'écart de 3 entre 7 et 4 n'est pas permis.}
\iqfollow{Et si vous ne disposez que de $O(k)$ de mémoire~?}{La file contient au plus $k$ positions, et les valeurs de $d$ ne doivent être conservées que pour celles-là, donc oui~: on les stocke dans la file, avec un meilleur courant.}
\end{solution}

\begin{solution}{iq:iv3:algorithms:79}
La somme pour \texttt{[2, 4, 1, 3]} vaut 17, et pour \texttt{[2, 2]} elle vaut 6. Le \Cref{lst:iv3:al:summins} compte, pour chaque valeur, le nombre de segments dont elle est le minimum, à l'aide d'une \index{pile monotone}pile monotone, en temps et mémoire $O(n)$.
\omcode{code/interviews-3/01-algorithms/python/iv_sum_mins.py}{1}{10}{Chaque valeur multipliée par le nombre de segments dont elle est le minimum.\protect\hfil\protect\break}\label{lst:iv3:al:summins}

Considérons une position $j$. Elle est le minimum d'un segment si le segment contient $j$ et rien de plus bas. Si la plus proche valeur inférieure à gauche est en position $\ell$, et la plus proche à droite en $r$, alors il y a $(j - \ell)$ choix pour le début et $(r - j)$ pour la fin, et la valeur compte donc $(j - \ell)(r - j)$ fois. Pour \texttt{[2, 4, 1, 3]}, la valeur 2 est le minimum de 2 segments, 4 de un, 1 de six et 3 de un~: $2 \times 2 + 4 \times 1 + 1 \times 6 + 3 \times 1 = 17$.

Les valeurs répétées sont le piège. Dans \texttt{[2, 2]}, le segment formé des deux a deux positions à égalité pour le minimum. Si les deux côtés s'arrêtent seulement à une valeur strictement plus basse, chaque 2 le revendique, et le total vaut $2 \times 2 + 2 \times 2 = 8$, en comptant ce segment deux fois. Si les deux côtés s'arrêtent à une valeur plus basse ou égale, personne ne le revendique, et le total vaut 4.

Il faut donc départager les égalités d'un seul côté~: on s'arrête à une valeur strictement plus basse à gauche, et à une valeur plus basse ou égale à droite, ou l'inverse. Alors un segment à égalité appartient à une seule position~: la dernière des égalités. La pile dépile ici sur \texttt{>=}, si bien que la limite droite de l'entrée dépilée est la première valeur inférieure ou égale, et que sa limite gauche est strictement plus basse. Chaque position est empilée et dépilée une seule fois, d'où $O(n)$.

La sentinelle supplémentaire vide la pile à la fin, et dans le test la méthode concorde avec une force brute sur tous les segments, sur des listes comportant beaucoup d'égalités.

\iqwrong{«~Somme sur les paires d'extrémités avec un minimum courant.~» C'est juste, mais en $O(n^2)$.}
\end{solution}

\begin{solution}{iq:iv3:algorithms:80}
La réponse est 9.75, obtenue à partir de \texttt{[10, -6, -5, 40]}. Le \Cref{lst:iv3:al:bestmean} est une \index{recherche dichotomique sur la réponse}recherche dichotomique sur la réponse, en temps $O(n \log(R/\epsilon))$ et en mémoire $O(1)$.
\omcode{code/interviews-3/01-algorithms/python/iv_best_mean.py}{1}{24}{Une moyenne de m est-elle atteignable~?}\label{lst:iv3:al:bestmean}

Une moyenne d'au moins $m$ est atteignable exactement lorsqu'il existe un segment de longueur au moins $k$ dont $\sum (x_i - m) \ge 0$. On retranche $m$, on prend les sommes préfixes, et en chaque extrémité $j$ on veut la plus petite somme préfixe prise au moins $k$ positions plus tôt~: si la somme préfixe courante est au-dessus de celle-ci, la réponse est oui. C'est un seul parcours. Et l'atteignabilité est monotone en $m$, ce qui permet à la recherche externe de fonctionner.

Un segment de longueur exactement $k$ ne suffit pas, car un segment plus long peut avoir une moyenne plus élevée. Sur \texttt{[9, 1, 9]} avec $k = 2$, chaque paire a pour moyenne 5, et les trois éléments ont $19/3$, soit environ 6.33. Il est vrai, en revanche, que les longueurs de $k$ à $2k - 1$ suffisent~: un segment plus long se découpe en deux morceaux d'au moins $k$ éléments, et l'un de ces morceaux a une moyenne au moins aussi élevée que l'ensemble.

Une méthode exacte simple consiste donc à essayer les longueurs de $k$ à $2k - 1$ avec des sommes glissantes, en $O(nk)$, ce que le test utilise comme seconde méthode, avec des fractions, et elle trouve la valeur exacte.

Cas limites~: avec $k = 1$, la réponse est la plus grande valeur isolée, et avec $k = n$, c'est la moyenne de tous les éléments. L'intervalle de recherche est sûr, car une moyenne ne peut pas sortir de l'intervalle des valeurs. Avec $k > n$, aucun segment n'existe, et la fonction renvoie \texttt{None} au lieu d'un nombre.

\iqwrong{«~L'algorithme de Kadane sur les valeurs.~» Il trouve la plus grande somme, et un long segment à grande somme peut avoir une faible moyenne, comme la liste entière ici~: somme 45, moyenne 7.5.}
\iqfollow{Et exactement~?}{Les moyennes de deux segments diffèrent d'au moins $1/n^2$ pour des valeurs entières, donc on s'arrête à cette tolérance et on retrouve la fraction à partir du meilleur segment. Il existe aussi une méthode linéaire utilisant l'enveloppe convexe des points des sommes préfixes.}
\end{solution}
